\documentclass[11pt,a4paper,openright,oneside]{book}
\usepackage{amsfonts, amsmath, amssymb,latexsym,amsthm, mathrsfs, enumerate}
\usepackage[utf8]{inputenc}
\usepackage[english]{babel}
\usepackage{epsfig}
\usepackage{csquotes}
\usepackage{biblatex}

\addbibresource{refs.bib}
\addbibresource{man-refs.bib}

\parskip=5pt
\parindent=15pt
\usepackage[margin=1.2in]{geometry}
\usepackage{graphicx}
\usepackage{listings}
\usepackage{tikz}
\usepackage{tikz-cd}
\usepackage{parskip}
\usepackage{algorithm}
\usepackage[noend]{algpseudocode}
\usepackage{hyperref}
\usepackage{cleveref}
\usepackage{color}
\hypersetup{
    linktoc=all,
    colorlinks=true,
    linkcolor=blue,  %choose some color if you want links to stand out
}

\setcounter{page}{0}



\numberwithin{equation}{section}
\newtheorem{defn0}{Definition}[chapter]
\newtheorem{prop0}[defn0]{Proposition}
\newtheorem{thm0}[defn0]{Theorem}
\newtheorem{lemma0}[defn0]{Lemma}
\newtheorem{corollary0}[defn0]{Corollary}
\newtheorem{example0}[defn0]{Example}
\newtheorem{remark0}[defn0]{Remark}
\newtheorem{conjecture0}[defn0]{Conjecture}

\newenvironment{definition}{ \begin{defn0}}{\end{defn0}}
\newenvironment{proposition}{\bigskip \begin{prop0}}{\end{prop0}}
\newenvironment{theorem}{\bigskip \begin{thm0}}{\end{thm0}}
\newenvironment{lemma}{\bigskip \begin{lemma0}}{\end{lemma0}}
\newenvironment{corollary}{\bigskip \begin{corollary0}}{\end{corollary0}}
\newenvironment{example}{ \begin{example0}\rm}{\end{example0}}
\newenvironment{remark}{ \begin{remark0}\rm}{\end{remark0}}
\newenvironment{conjecture}{\begin{conjecture0}}{\end{conjecture0}}

\newcommand{\defref}[1]{Definition~\ref{#1}}
\newcommand{\propref}[1]{Proposition~\ref{#1}}
\newcommand{\thmref}[1]{Theorem~\ref{#1}}
\newcommand{\lemref}[1]{Lemma~\ref{#1}}
\newcommand{\corref}[1]{Corollary~\ref{#1}}
\newcommand{\exref}[1]{Example~\ref{#1}}
\newcommand{\secref}[1]{Section~\ref{#1}}
\newcommand{\remref}[1]{Remark~\ref{#1}}
\newcommand{\conjref}[1]{Conjecture~\ref{#1}}
\newcommand{\figref}[1]{\cref{#1}}
\newcommand{\refeq}[1]{\cref{#1}}



\DeclareMathOperator{\vectorize}{vec}
\DeclareMathOperator{\rank}{rank}
\DeclareMathOperator{\unfolding}{unfold}
\DeclareMathOperator{\IN}{IN}
\DeclareMathOperator{\OUT}{OUT}
\DeclareMathOperator{\TNS}{TNS}
\DeclareMathOperator*{\argmax}{arg\,max}
\DeclareMathOperator*{\argmin}{arg\,min}
\DeclareMathOperator{\Tr}{Tr}
\DeclareMathOperator{\size}{Size}

% --------------------------------------------------
\usepackage{fancyhdr}

\lhead{}
\lfoot{}
\rhead{}
\cfoot{}
\rfoot{\thepage}

\begin{document}

\bibstyle{plain}

\thispagestyle{empty}

\begin{titlepage}
\begin{center}
\begin{figure}[htb]
\begin{center}
\includegraphics[width=6cm]{matematiquesinformatica-pos-rgb.png}
\end{center}
\end{figure}

\vspace*{1cm}
\textbf{\LARGE GRAU DE MATEM\`{A}TIQUES } \\
\vspace*{.5cm}
\textbf{\LARGE Treball final de grau} \\

\vspace*{1.5cm}
\rule{16cm}{0.1mm}\\
\begin{Huge}
\textbf{OPTIMAL LOW-RANK APPROXIMATION USING TENSOR NETWORK STRUCTURE SEARCH} \\
\end{Huge}
\rule{16cm}{0.1mm}\\

\vspace{1cm}

\begin{flushright}
\textbf{\LARGE Autor: Aran Roig}

\vspace*{2cm}

\renewcommand{\arraystretch}{1.5}
\begin{tabular}{ll}
\textbf{\Large Director:} & \textbf{\Large Dr. Nahuel Statuto} \\
\textbf{\Large Realitzat a:} & \textbf{\Large  Departament de Matemàtiques   } \\
 & \textbf{\Large i Informàtica} \\
\\
\textbf{\Large Barcelona,} & \textbf{\Large \today }
\end{tabular}

\end{flushright}

\end{center}



\end{titlepage}


\newpage
\pagenumbering{roman} 

\section*{Abstract}

Tensor network structure search has been interesting research topic since the raise on
complexity of deep learning models and quantum mechanics. This is an Undegraduate Thesis whose main goal is to give an 
automated search of an optimal tensor network structure for representing a given
tensor with some fixed error.

For these purpose we first give an introduction to tensors, tensor networks and then we present some examples
of well studied tensor networks, including Tucker decomposition, Tensor Train decomposition, Tensor Ring decomposition
and Fully Connected Tensor Network decomposition. 

Then with some practical experiments we demonstrate that it is possible to find more optimized structures
without significant losses on performance and accuracy, and finally we will present
an algorithm for finding these optimized structures.

\section*{Resum}

La recerca de l'estructura òptima de xarxes de tensors ha estat un tema d'interès des de l'augment en la complexitat dels models d'aprenentatge profund i de la mecànica quàntica. 
Aquest és un treball de final de grau que té com a objectiu principal oferir una cerca automatitzada 
d’una estructura òptima de xarxa de tensors per representar un tensor donat amb un error fixat.

Per aconseguir aquest propòsit, primer oferim una introducció als tensors, a les xarxes de tensors, i tot seguit presentem alguns exemples de xarxes de tensors ben estudiades, 
incloent-hi la descomposició de Tucker, la descomposició en tren de tensors (Tensor Train), la descomposició en anell
de tensors (Tensor Ring) i la descomposició de xarxa de tensors totalment connectada.

A continuació, amb alguns experiments pràctics, demostrem que és possible trobar estructures més optimitzades sense pèrdues significatives 
en el rendiment i la precissió de la representació, i finalment presentem un algoritme per trobar aquestes estructures optimitzades.

% TODO: Omplir això
% 15A69 - Multilinear Algebra, Tensor Calculus, Graph Theory, no se
{\let\thefootnote\relax\footnote{2020 Mathematics Subject Classification. 11G05, 11G10, 14G10}}



\newpage 


\section*{Agra\"{\i}ments}

Vull agrair a ... 
\newpage

{\hypersetup{linkcolor=black}
\tableofcontents
}

\newpage

\pagenumbering{arabic} 
\setcounter{page}{1}
\chapter{Introduction}

Neural networks are a set of dependant non-linear functions which each individual function corresponds to 
a neuron (or perceptron). In a fully connected neural network, each neuron applies a linear transformation
to the input vector through a weight of matrices.

The output of a neuron $j$ in some layer is given by the formula:

$$ y_j(x) = f\left(\sum_{i=1}^N w_{ji} x_i + b_{j}\right) $$ 

Where $w_{ji}$ is the weight assigned to the connection between the neuron $i$ of the previous layer with the neuron
$j$ of the current layer, $f$ is an activation function and $b_j$ is some bias of the neuron $j$.

We can represent the weights of a fully connected neural network as a matrix, that we call the \textbf{weight matrix}
since the recent rise on complexity of neural networks, weight matrices can be very large and there is a lot of
interest on compressing weight matrices without sacrificing accuracy or performance.

Since we can reshape the weight matrix into a tensor, in this thesis, we will aim to compress a tensor
using tensor networks, a concept that originates from the study of many-body quantum systems \cite{orusTensorNetworksComplex2019}
and recently has gained traction on machine learning.

The objective of this thesis is given a tensor $T$, finding the tensor network structure that
is represented by tensors of smaller rank (called core tensors) that, when the whole
network is evaluated, represent $T$ well enough given a fixed error $\epsilon$. We will also minimize the computational cost 
when it comes to contracting the whole network for retrieving the representation of the original tensor.

We will see that finding the best structure is an integer programming problem, that is $NP$-hard. We will deduce
and give an algorithm that finds a locally best structure by using program synthesis, and finally, we will compare this 
structure with some other more studied cases of tensor networks, as Tensor Train networks, Tensor Ring networks and
fully connected tensor networks

%\section{Objectives}

\section{Thesis structure}

First we will present some preliminaries about tensor algebra. Then, we will introduce the
diagrammatic notation made by Roger Penrose in the earlies 1970. We will
use it for representing both matrix product states and tensor networs.

Then, we will give a formal introduction in tensor networks, based in most part from \cite{yeTensorNetworkRanks2019}.
We will introduce also well studied tensor networks such as the Tucker decomposition, the tensor train decomposition, the tensor ring decomposition,
the fully connected tensor network decomposition.

We will make special emphasis on tree tensor networks, since they have some nice properties that we can later use
for splitting tensor networks. We will also describe the alternating least squares algorithm applied to tensor networks
for finding approximated states for any tensor network.

Finally, we will explain some algorithms for finding optimal tensor network structures. We will give an algorithm for
finding a general tensor network structure based following \cite{liPermutationSearchTensor2022} and we will also describe an algorithm for finding an 
optimal tensor network tree structure using program synthesis following mainly from the paper \cite{guoTensorNetworkStructure2025}



\chapter{Preliminaries}


\section{Graph Theory}

\begin{definition}[Graph]
    \normalfont{\cite{wilsonIntroductionGraphTheory2009}}
    A \textbf{graph} $G$ is defined as a tuple $G = (V,E)$ where $V = V(G)$ is a set of elements called \textbf{vertices} and 
    $E = E(G) \subset \{\{u, v\} : u, v \in V\}$ is a set of elements called \textbf{edges}.
\end{definition}

\begin{definition}[Directed Graph]
    A \textbf{directed graph} $D$ is a tuple $D = (V, \bar{E})$ where $V = V(D)$ are its vertices and
    $\bar{E} = \bar{E}(D) \subset \{(u,v) : u,v \in V\}$
\end{definition}

\begin{definition} Given a directed graph $G=(V,\bar{E})$ and a vertex $i \in V$ we define
    $$\IN(i) = \{j \in V : (j, i) \in \bar{E}\} \qquad \OUT(i) = \{j \in V : (i, j) \in \bar{E} \}$$
\end{definition}
\iffalse
\begin{definition}[Graph isomorphisms]
    We say that two graphs $G, H$ are \textbf{isomorphic} if 
\end{definition}


\begin{definition}[Outer product]
    We define the \textbf{outer product} of $n$ vectors 
    $a_1 \in \mathbb{R}^{N_1}, a_2 \in \mathbb{R}^{N_2}, \dots, a_n \in \mathbb{R}^{N_n}$
    as the $n$th-order tensor $\mathcal{T} \in \mathbb{R}^{N_1 \times N_2 \times \dots \times N_n}$ with its entries being
    $$\mathcal{T}(i_1, \dots, i_n) = a_1(i_1) a_2(i_2) \cdots a_n(i_n)$$
\end{definition}

\begin{example}
The outer product of two vectors $a \in \mathbb{R}^I$ and $b \in \mathbb{R}^J$ is denoted by $a \circ b$ and it results
as a matrix $M = a \circ b \in \mathbb{R}^{I \times J}$ with its entries defined as $M_{ij} = a_i b_j$.
\end{example}


\begin{definition}[Inner product]
The inner product of two tensors $\mathcal{X}, \mathcal{Y} \in \mathbb{R}^{N_1 \times \cdots \times N_n}$ is defined by
$$\langle \mathcal{X},\mathcal{Y} \rangle = \sum_{i_1, \dots, i_N}^{N_1, \dots, N_n} \mathcal{X}_{i_1, \dots, i_n} \mathcal{Y}_{i_1, \dots, i_n} = 
\vectorize(\mathcal{X})^T \vectorize(\mathcal{Y}) = \langle \vectorize(\mathcal{X}), \vectorize(\mathcal{Y}) \rangle$$
\end{definition}


\begin{definition}[Kronecker product]
    \normalfont{\cite{panagakisTensorMethodsComputer2021}} Given two matrices $A \in \mathbb{R}^{N_1 \times N_2}$ and $B \in \mathbb{R}^{M_1 \times M_2}$,
    their kronecker product is defined as the matrix $A \otimes B \in \mathbb{R}^{N_1 \cdot M_1 \times N_2 \cdot M_2}$ with
    $$A \otimes B = \begin{bmatrix}
        a_{11}B & \cdots & a_{1N_2}B \\ 
        \vdots & \ddots & \vdots \\
        a_{N_1 1}B & \cdots & a_{N_1 N_2}B \\
    \end{bmatrix}$$
\end{definition}

\begin{definition}[Khatri-Rao product]
    \normalfont{\cite{panagakisTensorMethodsComputer2021}} Given two matrices $A \in \mathbb{R}^{N \times R}$ and $B \in \mathbb{R}^{M \times R}$ their
Khatri-Rao, also known as column-wise Kronecker product is defined as $A \odot B \in \mathbb{R}^{N \cdot M \times R}$
    $$ A \odot B = \begin{bmatrix} A_{:,1} \otimes B_{:,1} & A_{:,2} \otimes B_{:,2} & \cdots & A_{:,R} \otimes B_{:,R}  \end{bmatrix}$$
        

\end{definition}

\begin{definition}[Frobenius norm]
The \textbf{Frobenius norm} of a tensor $\mathcal{T} \in \mathbb{R}^{N_1 \times \cdots \times N_n}$ is given by
$$\|\mathcal{T}\|_F = \sqrt{\langle \mathcal{T}, \mathcal{T} \rangle} = \sqrt{\sum_{i_1, \dots, i_n}^{N_1, \dots, N_n}
\mathcal{T}_{i_1 \dots i_n}^2}$$
\end{definition}

\begin{definition}[Tensor transposition]
    \normalfont{\cite{zhengFullyConnectedTensorNetwork2021}} Let $\mathcal{T} \in \mathbb{R}^{N_1 \times \cdots \times N_n}$
    be an $n$th-order tensor and $p$ a permutation of the vector $(1, 2, \dots, n)$. We define the \textbf{vector $p$ based tensor
    transposition of $\mathcal{T}$} as the tensor $\overrightarrow{\mathcal{T}_p} \in \mathbb{R}^{N_{p_1} \times \cdots \times N_{p_n}}$ with its entries defined as follows:
    $$\overrightarrow{\mathcal{T}_p}(i_1, i_2, \cdots, i_n) = (i_{p_1}, i_{p_2}, \cdots, i_{p_n})$$
\end{definition}

\begin{definition}[Tensor contraction]
    \normalfont{\cite{zhengFullyConnectedTensorNetwork2021}} Suppose that $p$ and $q$ are reorderings of the vectors
    $(1,2,\dots,n)$ and $(1,2,\dots,m)$ respectively, and let ${\mathcal{X} \in \mathbb{R}^{N_1 \times \cdots \times N_n}}$ 
    and $\mathcal{Y} \in \mathbb{R}^{M_1 \times \cdots \times M_m}$ two tensors with $N_{p_i} = M_{q_i}$ for all $i = 1,2,\dots,d$
    with $d \leqslant \min{(n, m)}$. We define the tensor contraction along the $p_{1:d}$-modes of $\mathcal{X}$ and the $q_{1:d}$-modes
    of $\mathcal{Y}$ as the tensor $\mathcal{Z}$ of order $n + m - 2d$
$$\mathcal{Z} = \mathcal{X} \times_{p_{1:d}}^{q_{1:d}} \mathcal{Y} \in \mathbb{R}^{N_{p_{d+1}} \times \cdots \times N_{p_{n}} \times N_{q_{d+1}} \times \cdots \times N_{q_m}}$$
whose elements are defined by:
$$\mathcal{Z}(i_{p_{d+1}}, \cdots, i_{p_n}, j_{q_{d+1}}, \cdots, j_{q_m}) = $$$$ \sum_{i_{p_1} = 1}^{N_1} \sum_{i_{p_2} = 1}^{N_2} \cdots \sum_{i_{p_d} = 1}^{N_d}
\overrightarrow{\mathcal{X}_p}(i_{p_1}, \cdots, i_{p_d}, i_{p_{d+1}}, \cdots, i_{p_n}) \overrightarrow{\mathcal{Y}_q}(i_{p_1}, \cdots, i_{p_d}, j_{q_{d+1}}, \cdots, j_{q_m})$$
\end{definition}



\fi




In this chapter we will introduce the basics of tensor algebra and we will present some results that will be useful in the
following chapters

\section{Basic tensor definitions}

We will denote $\mathbb{V}_1, \dots, \mathbb{V}_n$ as finite vector spaces over a field $\mathbb{K}$ ($\mathbb{C}$ if unspecified) of dimension $\dim{\mathbb{V}_i} = N_i \; \forall i = 1, \dots, n$.
We denote $\mathcal{B}_1, \dots, \mathcal{B}_n$ with ${\mathcal{B}_i = \{e_1^{(i)}, \dots, e_{N_i}^{(i)}\}, i = 1, \dots, n}$ the canonical basis of $\mathbb{V}_1, \dots, \mathbb{V}_n$

\begin{definition}[Tensor]
    A \textbf{tensor} is a multilineal map ${T: \mathbb{V}_1 \times \dots \times \mathbb{V}_n \rightarrow \mathbb{K}}$
\end{definition}

\begin{definition}[Order of a tensor] Given a tensor $T \in \mathbb{V}_1 \times \cdots \times \mathbb{V}_n \rightarrow \mathbb{K}$ We define the \textbf{order} of the tensor $T$ as $n$.
\end{definition}

We will allow that some $\mathbb{V}_i$ are dual spaces of some other vector spaces $\mathbb{W}_i$, i.e $\mathbb{V}_i = \mathbb{W}_i^*$.
\begin{definition}
    Let $T \in \mathbb{V}_1 \times \cdots \times \mathbb{V}_p \times \mathbb{W}_1^* \times \cdots \times \mathbb{W}_q^* \rightarrow \mathbb{K}$.
    We say that the tensor $T$ is $p$-times covariant and $q$-times contravariant.
\end{definition}

\begin{definition}[Tensor product] Let $L$ be the vector space generated by the base ${V_1 \times \dots \times V_n}$, i.e the
    set of linear combinations of the elements $(v_1, \dots, v_n), v_i \in \mathbb{V}_i$. Let $\mathcal{R}$ be the linear subspace of
    $L$ generated by the relation $R$ defined by:
    $$(v_1, \dots, \alpha v_i, \dots, v_n) \sim \alpha(v_1, \dots, v_n) \; \forall i = 1, \dots, n, \forall \alpha \in \mathbb{K}$$
    $$(v_1, \dots, v_i + u_i, \dots, v_n) \sim (v_1, \dots, v_i \dots, v_n) + (v_1, \dots, u_i, \dots, v_n) \; \forall i = 1, \dots, n$$
    The tensor product $\mathbb{V}_1 \otimes \dots \otimes \mathbb{V}_n$ is defined as the quotient $L / \mathcal{R}$ and its called 
    \textbf{tensor product space}. The image of
    $(v_1, \dots, v_n)$ by the quotient is denoted by $v_1 \otimes \dots \otimes v_n$
\end{definition}

The following theorem gives us a correspondance between each tensor ${T: \mathbb{V}_1 \times \dots \times \mathbb{V}_n \rightarrow \mathbb{K}}$
and each element of $\mathbb{V}_1 \otimes \dots \otimes \mathbb{V}_n$:

\begin{theorem}[Universal property of the tensor product]
    \normalfont{\cite{romanTensorProducts2008}}
    The tensor product of two vector spaces $\mathbb{V} \otimes \mathbb{W}$
    for every bilinear map $h: \mathbb{V} \times \mathbb{W} \rightarrow X$ there exists an unique bilinear map $\tilde{h}: \mathbb{V} \otimes
    \mathbb{W} \rightarrow X$ such that the following diagram commutes:

    \centering
% https://q.uiver.app/#q=WzAsMyxbMCwwLCJWXFx0aW1lcyBXIl0sWzEsMCwiViBcXG90aW1lcyBXIl0sWzEsMSwiWCJdLFswLDEsIlxcdmFycGhpIl0sWzAsMiwiaCIsMl0sWzEsMiwiXFx0aWxkZSBoIl1d
\begin{tikzcd}
	{\mathbb{V}\times \mathbb{W}} & {\mathbb{V} \otimes \mathbb{W}} \\
	& X
	\arrow["\varphi", from=1-1, to=1-2]
	\arrow["h"', from=1-1, to=2-2]
	\arrow["{\tilde h}", from=1-2, to=2-2]
\end{tikzcd}

\end{theorem}

We can construct a basis for $\mathbb{V}_1 \otimes \dots \otimes \mathbb{V}_n$. We define
$$\mathcal{B}_{\otimes} = \{e_{i_1}^{(1)} \otimes \cdots \otimes e_{i_n}^{(n)} : 1 \leqslant i_j \leqslant N_j, 1 \leqslant j \leqslant n\}$$
Where $\{e_1^{(i)}, e_2^{(i)}, \dots, e_{N_i}^{(i)}\}$ is the canonical basis for $\mathbb{V}_i$.
Constructed this way, $\mathcal{B}_\otimes$ is a (canonical) basis of ${\mathbb{V}_1 \otimes \dots \otimes \mathbb{V}_n}$.

\begin{remark} \label{rem:tenbase} The dimension
    of ${\mathbb{V}_1 \otimes \dots \otimes \mathbb{V}_n}$ is ${\dim{\mathbb{V}_1} \cdot \dim{\mathbb{V}_2} \cdots \dim{\mathbb{V}_n}}$ and its elements can be expressed as
    \begin{equation} \label{eq:base-representation}
T = \sum_{s_1, \dots, s_n}^{N_1, \dots, N_n} T_{s_1, \dots, s_n} \cdot  e_{s_1}^{(1)} \otimes \cdots \otimes e_{s_n}^{(n)}
\end{equation}
We will define the size of the tensor $T$ as $\size(T) = \dim \mathbb{V}_1 \cdot \dim \mathbb{V}_2 \cdots \dim \mathbb{V}_n$
\end{remark}


% (POSAR COSES DEL KROENKER PRODUCT, PRESENTARLO)
\begin{definition}[Tensor product] Given two tensors $T \in \mathbb{V}_1 \otimes \cdots \otimes \mathbb{V}_n$ and ${U \in \mathbb{W}_1 \otimes \cdots \otimes
    \mathbb{W}_m}$ with $\dim \mathbb{V}_i = N_i$, $i = 1, \dots, n$ and $\dim \mathbb{W}_j = M_j$, $j = 1, \dots, m$ 
    and canonical basis $\{v_1^{(i)}, \dots, v_{N_i}^{(i)}\}$ for each $\mathbb{V}_i$ and $\{w_1^{(j)}, \dots, w_{M_j}^{(j)}\}$ for each $\mathbb{W}_j$
    we define the tensor product
    $T \otimes U$ as
    $$ T \otimes U = \sum_{i_1, \dots, i_n}^{N_1, \dots, N_n} \sum_{j_1, \dots, j_m}^{M_1, \dots, M_m} T_{i_1, \dots, i_n} U_{j_1, \dots, j_m} \cdot
    v_{i_1}^{(1)} \otimes \cdots \otimes v_{i_n}^{(n)} \otimes w_{j_1}^{(1)} \otimes \cdots \otimes w_{j_m}^{(m)}$$
    Note that $T \otimes U \in \mathbb{V}_1 \otimes \cdots \otimes \mathbb{V}_n \otimes \mathbb{W}_1 \otimes \cdots \otimes \mathbb{W}_m$

\end{definition}

From \remref{rem:tenbase}, a tensor $T \in \mathbb{V}_1 \otimes \cdots \otimes \mathbb{V}_n$ can be identified as a "$n$-dimensional array" $\mathcal{T}$, i.e,
a discrete function
$$\begin{align}
    \mathcal{T}: \prod_{i=1}^n \{1, \dots, N_i\} & \longrightarrow \mathbb{K} \\
    T(i_1, \dots, i_n) & \longmapsto T_{i_1, \dots, i_n}
\end{align}$$

From now on, we will write the discrete function $\mathcal{T}$ as
an element of ${\mathbb{K}^{N_1 \times \cdots \times N_n}}$ \cite{yokotaVeryBasicsTensors2024}.
of $T$ as $\mathcal{T}$. Now we will define some definitions from the underlying definition of the tensor viewed as a discrete function:

Now we will define some definitions that will help us establish a formalized method of mapping $n$-dimensional arrays with
vectors and matrices:

\begin{definition}[Linearization]
    Fixed $N_1, \dots, N_n \in \mathbb{N}$, given $i_1, \dots, i_n \in \mathbb{N}$ such that $1 \leqslant i_1, \leqslant N_1, \dots, 1 \leqslant i_n \leqslant N_n$,
    we define the Linearization of the indices $i_1, \dots, i_n$ as the mapping ${\prod_{i=1}^n \{1, \dots, N_i\} \rightarrow \{1, \dots, \prod_{i=1}^n N_i\}}$:
    $$\overline {i_1, i_2, \dots, i_n} = \sum_{j=2}^{n} \left( (i_j - 1) \prod_{k=1}^j N_k \right) + i_1$$
\end{definition}

\begin{lemma}
    The linearization mapping is bijective
\end{lemma}

\begin{definition}[Vectorization]
    Given a tensor $\mathcal{T} \in \mathbb{K}^{N_1 \times \cdots \times N_n}$, we define the \textbf{vectorization} of $\mathcal{T}$
    as the first order tensor $\mathcal{V} \in \mathbb{K}^{N_1 N_2 \cdots N_n}$ defined by:
    $$\mathcal{V}(\overline{i_1 i_2 \dots i_n}) = \mathcal{T}(i_1, i_2, \dots, i_n)$$
    We will write $\mathcal{V} = \vectorize{\mathcal{T}}$
\end{definition}


\begin{definition}[Tensor unfolding]
Let $\mathcal{T} \in \mathbb{R}^{N_1 \times \cdots \times N_n}$, $n \geqslant 2$, $1 \leqslant d \leqslant n$ and $p$ a permutation of the vector $(1,2,\dots, n)$. We define the
\textbf{generalized unfolding} of the tensor $\mathcal{T}$ as the $2$nd-order tensor 
$\mathcal{U} \in \mathbb{R}^{\prod_{i=1}^d N_{p_i} \times \prod_{i=d+1}^n N_{p_i}}$:
$$ \mathcal{U} (\overline{i_{p_1}, \dots, i_{p_d}}, \overline{i_{p_{d+1}}, \dots, i_{p_n}}) = \mathcal{T}(i_1, \dots, i_n)$$

We will write $\mathcal{U} = \unfolding{(\mathcal{T}, (p_1, \dots, p_d), (p_{d+1}, \dots, p_n))}$. 
We also define $\unfolding_d T := \unfolding{(\mathcal{T}, (1, \dots, d), (d+1, \dots, n))}$
\end{definition}

\begin{figure}[h]
    \centering
    \begin{tikzpicture}
        \pgfmathsetmacro{\cubex}{2}
        \pgfmathsetmacro{\cubey}{2}
        \pgfmathsetmacro{\cubez}{2}
        \pgfmathsetmacro{\coordx}{0}
        \pgfmathsetmacro{\coordy}{0}
        \pgfmathsetmacro{\coordz}{0}
        \draw[fill=lightgray] (\coordx,\coordy,\coordz) -- ++(-\cubex,0,0) -- ++(0,-\cubey,0) -- ++(\cubex,0,0) -- cycle;
        \draw[fill=lightgray] (\coordx,\coordy,\coordz) -- ++(0,0,-\cubez) -- ++(0,-\cubey,0) -- ++(0,0,\cubez) -- cycle;
        \draw[fill=lightgray] (\coordx,\coordy,\coordz) -- ++(-\cubex,0,0) -- ++(0,0,-\cubez) -- ++(\cubex,0,0) -- cycle;
        

        \node(j3) at (-1, -1) {$T$};
        

    \end{tikzpicture}
    \label{fig:mode_unfolding}
    \caption{Tensor unfolding}
\end{figure}

\begin{definition}[Tensor slices]
%Consider a matrix $A = [\mathbf{a}_1, \mathbf{a}_2, \dots, \mathbf{a}_J] = [\mathbf{\tilde a}_1, \mathbf{\tilde a}_2, \dots, \mathbf{\tilde a}_I]^T \in \mathbb{R}^{I \times J}$. We define the slices of a matrix
%as $$A(:, j) = \mathbf{a}_j \in \mathbb{R}^I \qquad A(i, :) = \mathbf{\tilde{a}}_j \in \mathbb{R}^J$$

    Let $\mathcal{T} \in \mathbb{K}^{N_1\times \cdots \times N_n}$. Consider $S \subset \{1, \dots, n\}$ a subset of modes (dimensions)
    of $\mathcal{T}$. Let $\mathbf{i}_S = (i_k)_{k \in S}$. We define the \textbf{slice of $\mathcal{T}$} as the tensor $\mathcal{T}_{\mathbf{i}_S}$
    of order $n - \#S$
    $$\begin{align}
        \mathcal{T}_{\mathbf{i}_S} : \prod_{k \not\in S} \{1, \dots, N_k\} & \longrightarrow \mathbb{K} \\
        \mathcal{T}_{\mathbf{i}_S}(i_{j_1}, \dots, i_{j_m}) & \longmapsto \mathcal{T}(i_1, \dots, i_n)
    \end{align}$$
    Where $\{j_1, \dots, j_m\} = \{1, \dots, n\} \setminus S$. 

    Sometimes we will also implicitly specify $\mathbf{i}_S$ by writing $\mathcal{T}(a_1, \dots, a_n)$ and replacing $a_j$ with $i_j$ if $j \in S$ and "$:$" otherwise

    We will denote $\mathcal{T}^{(m)} := \mathcal{T}(:, \dots, :, i_m, :, \dots, :)$ and $\mathcal{T}^{(\neq m)} := \mathcal{T}(i_1, \dots, i_{m-1}, :, i_{m+1}, \dots, i_n)$
\end{definition}

\begin{figure}[h]
    \centering
    \begin{tikzpicture}
        \pgfmathsetmacro{\cubex}{2}
        \pgfmathsetmacro{\cubey}{2}
        \pgfmathsetmacro{\cubez}{2}
        \pgfmathsetmacro{\coordx}{0}
        \pgfmathsetmacro{\coordy}{0}
        \pgfmathsetmacro{\coordz}{0}
        \draw[fill=lightgray] (\coordx,\coordy,\coordz) -- ++(-\cubex,0,0) -- ++(0,-\cubey,0) -- ++(\cubex,0,0) -- cycle;
        \draw[fill=lightgray] (\coordx,\coordy,\coordz) -- ++(0,0,-\cubez) -- ++(0,-\cubey,0) -- ++(0,0,\cubez) -- cycle;
        \draw[fill=lightgray] (\coordx,\coordy,\coordz) -- ++(-\cubex,0,0) -- ++(0,0,-\cubez) -- ++(\cubex,0,0) -- cycle;
        

        \node(j3) at (-1, -1) {$T$};
        

    \end{tikzpicture}
    \label{fig:slicing}
    \caption{Tensor slices}
\end{figure}


\begin{example}
    Consider $\mathcal{M} \in \mathbb{K}^{N_1 \times N_2}$ a second order tensor. We can see this tensor as a
    bidimensional array (matrix). The slice $\mathcal{M}(i, :)$ results in the $i$-th row of $\mathcal{M}$.
    The slice $\mathcal{M}(:, j)$ results in the $j$-th column of $\mathcal{M}$
\end{example}

\begin{example}
    Consider a $4$th-order tensor $\mathcal{T} \in \mathbb{K}^{N_1 \times N_2 \times N_3 \times N_4}$. 
    Fixed $i_2, i_3$, the tensor slice of $2$th-order $\mathcal{T} = \mathcal{A}(:, i_2, i_3, :) \in \mathbb{K}^{N_1 \times N_4}$ 
    with its entries defined by ${\mathcal{T}(i_1, i_4) = \mathcal{A}(i_1, i_2, i_3, i_4)}$
\end{example}

\begin{definition}[Rank of a tensor] 
    We say that a tensor is of rank $r$ and we write $\rank{T} = r$
    with $r \in \mathbb{N}$ being the minimum value such that we can write $T$ as
$$T= \sum_{p=1}^r v_p^{(1)} \otimes \cdots \otimes v_p^{(n)}$$
where $v_1^{(i)}, \dots, v_r^{(i)} \in \mathbb{V}_i, i = 1, \dots, n$
\end{definition}

One can easily see that $\rank{T} \leqslant \prod_{i=1}^n N_i$.
Unlike matrices, determining the rank of a tensor is an NP-hard problem. \cite{hillarMostTensorProblems2013}. Even
finding the maximum rank, (i.e determining $\max_{T \in \mathbb{V}_1 \otimes \cdots \otimes \mathbb{V}_n} \rank{T}$) this is still an unresolved problem.
We will now present an slightly better upper bound for the tensor rank:
\begin{proposition}
    \begin{equation} \label{eq:rank-dimensional-bound}
        \rank{T} \leqslant \left\lfloor \frac{\prod_{i=1}^n N_i}{\sum_{i=1}^n N_i} \right\rfloor
    \end{equation}

\begin{proof}
    Let $r = \rank{T}$. We can write $T = \sum_{p=1}^r v_p^{(1)} \otimes \cdots \otimes v_p^{(n)}$. Now, each term of this
    sum has $\sum_{i=1}^n N_i$ adjustable parameters, since each $v_p^{(i)}$ is a vector of $\mathbb{V}_i$ with its dimension being $N_i$.
    So, in total we will have $r \sum_{i=1}^n N_i$ adjustable parameters in our decomposition. Since our tensor $T$ is completly
    determined by $\prod_{i=1}^n N_i$ parameters, we can impose that $r \sum_{i=1}^n N_i \leqslant \prod_{i=1}^n N_i$ 
\end{proof}
\end{proposition}

\chapter{Tensor networks}

\section{Penrose Notation}

The Penrose an intuitive graphical language to represent tensor contractions that
dates back from at least the early 1970s \cite{rogerPenroseApplications}

Given an $n$th-order tensor $\mathcal{T} \in \mathbb{K}^{N_1 \times \dots \times N_n}$ the way we draw it using the
Penrose notation is as a circle with as many edges as the order of the tensor, as seen in \figref{fig:tens}

\begin{figure}[h]
\centering
\tikz {
    \node(T)[draw, shape=circle, minimum size=0.8cm] at (0,0) {$\mathcal{T}$};
    \node(i) at (-1, 0) {$i_1$};
    \node(i2) at (-0.7, 0.8) {$i_2$};
    \node(id) at (0.7, 0.8) {$\dots$};
    \node(in) at (1, 0) {$i_n$};
    \draw (T) -- (i);
    \draw (T) -- (i2);
    \draw (T) -- (id);
    \draw (T) -- (in);
}
\caption{
    Representation of a tensor $\mathcal{T} \in \mathbb{K}^{N_1 \times \cdots \times N_n}$ using the Penrose notation
}
\label{fig:tens}
\end{figure}

A lot of times we will not write explicitly the name of the indexes, since we only care about their order. The order
of the indexes will be determined by their orientation respect to the circle: to get the order we will start from the left and then
following a clockwise rotation. The order in which we encounter the edges will be the order of the indexes. For example,
in \figref{fig:tens}, the order would be $i_1, i_2, \dots, i_n$


\section{Tensor contraction}

We would want to generalize the concept of vector dot product, or matrix product or the trace operation into tensors.
The tensor contraction gives us a generalization of all of the above. The concept of the tensor contraction 
basically arises from the canonical pairing of a vector space and its dual. Since tensors are elments
of $T \in \mathbb{V}_1 \otimes \cdots \otimes \mathbb{V}_n \otimes \mathbb{W}_1 \otimes \cdots \otimes \mathbb{W}_m$ with for each
$\mathbb{W}_j$ being the dual space of some space $\mathbb{V}_i$, one can contract these tensor by pairing these two
spaces.

Formally, the contraction of two tensors $\mathcal{X}, \mathcal{Y}$ is the image of a fixed contraction $\mathcal{C}_k^l$ map where $k$ is an index (or edge) of the first tensor $\mathcal{X}$ and $l$
is another index of the second tensor $\mathcal{Y}$:

\begin{definition}
    \normalfont{\cite{dodsonTensorsMultilinearForms1991}}
    Let $T \in \left( \bigotimes_{i=1}^{k-1} \mathbb{V}_i \otimes \mathbb{V}_k \otimes \bigotimes_{i=k+1}^{p} \mathbb{V}_i \right) \otimes 
    \left( \bigotimes_{i=1}^{l-1} \mathbb{W}_{i}^* \otimes \mathbb{W}_l^* \otimes 
    \bigotimes_{i=l+1}^q \mathbb{W}_i^* \right)$ With $\mathbb{V}_k = \mathbb{W}_l$. Consider the mapping
    $$
    \mathcal{C}_k^l: 
    \left(\bigotimes_{i=1}^{p} \mathbb{V}_i \right) \otimes \left( \bigotimes_{i=1}^{q} \mathbb{W}_i^* \right) \longrightarrow 
    \left( \bigotimes_{i=1}^{k-1} \mathbb{V}_i \otimes \bigotimes_{i=k+1}^p \mathbb{V}_i \right) \otimes \left( \bigotimes_{i=1}^{l-1} \mathbb{W}^*_i \otimes \bigotimes_{i=l+1}^q \mathbb{W}^*_i \right)
$$
$$
\mathcal{C}_k^l \left( \bigotimes\nolimits_{i=1}^p v_i \otimes \bigotimes\nolimits_{i=1}^q f_q \right) 
= \left(\bigotimes\nolimits_{i=1}^{k-1} v_i \otimes \bigotimes\nolimits_{i=k+1}^p v_{i} \otimes
    \bigotimes\nolimits_{i=1}^{l-1} f_i \otimes \bigotimes_{i=l+1}^q f_i \right) f_l(v_k)$$
$\mathcal{C}_k^l$ is defined as the tensor contraction mapping of $\mathbb{V}_1 \otimes \cdots \otimes \mathbb{V}_p \otimes \mathbb{V}_{p+1}^* \otimes
\cdots \otimes \mathbb{V}_q^*$ by the indices $(k,l)$. We call $\mathcal{C}_k^l(T)$ the contraction of $T$ by indices $(k,l)$. In other words, we apply
the dual space $\mathbb{V}_k^*$ with $\mathbb{V}_k$.

Now we will introduce the definition of the contraction between two tensors by defining a contraction over the tensor product
of the two original tensors.
\end{definition}

\begin{definition}

    Let $X \in \left( \bigotimes_{i=1}^{k-1} \mathbb{V}_i \otimes \mathbb{V}_k \otimes \bigotimes_{i=k+1}^{p} \mathbb{V}_i \right)$ and
    $Y \in \left( \bigotimes_{i=1}^{l-1} \mathbb{W}_{i} \otimes \mathbb{W}_l^* \otimes 
    \bigotimes_{i=l+1}^q \mathbb{W}_i \right)$ with $\mathbb{V}_k = \mathbb{W}_l$. We define the contraction between $X$ and $Y$ by the
    indices $(k,l)$ as $$X \times^l_k Y := \mathcal{C}_k^l(X \otimes Y)$$
\end{definition}

We will represent the contraction between two tensors 
as their representation in the Penrose notation with the edges that represent the indexes that are contracting by joining them, as seen in \figref{fig:tencon}.


\begin{figure}
    \begin{center}
    \hfill
\begin{minipage}{0.2\textwidth}
    \begin{tikzpicture}
    \node(X)[draw, shape=circle, minimum size=0.8cm] at (0,0) {$\mathcal{X}$};
    \node(Y)[draw, shape=circle, minimum size=0.8cm] at (2,0) {$\mathcal{Y}$};
    \node(i1) at (-1, 0) {$i_1$};
    \node(i2) at (0, 1) {$i_2$};
    \node(j2) at (2, 1) {$j_2$};
    \node(j3) at (3, 0) {$j_3$};
    \draw (X) -- (i1);
    \draw (X) -- (i2);
    \draw (Y) -- (j2);
    \draw (Y) -- (j3);
    \draw (X) -- (Y) 
        node[above, pos=0.8] {$j_1$}
        node[above, pos=0.2] {$i_3$};
\end{tikzpicture}
\end{minipage}
\hfill
\begin{minipage}{0.225\textwidth}
\begin{tikzpicture}
    \node(1) at(-1.5, -0.5);
    \node(2) at(3, -0.5);
    \draw[-{Stealth[length=12pt, inset=2pt]}] (1) -- (2)
        node[above, pos=0.45] {$\mathcal{C}_{i_3}^{j_1}$};
\end{tikzpicture}
\end{minipage}
\hfill
\begin{minipage}{0.225\textwidth}
\begin{tikzpicture}
    \node(Z)[draw, shape=circle, minimum size=0.8cm] at (0,0) {$\mathcal{Z}$};
    \node(i1) at (-1, 0) {$i_1$};
    \node(i2) at (-0.4, 1) {$i_2$};
    \node(j2) at (0.4, 1) {$j_2$};
    \node(j3) at (1, 0) {$j_3$};
    \draw (Z) -- (i1);
    \draw (Z) -- (i2);
    \draw (Z) -- (j2);
    \draw (Z) -- (j3);
\end{tikzpicture}
\end{minipage}
\end{center}

\caption{
    Representation in the Penrose notation of the contraction between two tensors $\mathcal{X} \in \mathbb{K}^{N_1 \times N_2 \times N_3}, \mathcal{Y} \in \mathbb{K}^{M_1 \times M_2 \times M_3}$
    by their indices $i_3$ and $j_1$ with $N_1, N_2, N_3, M_1, M_2, M_3 \in \mathbb{N}$ and $N_3 = M_1$
}
\label{fig:tencon}
\end{figure}

Fixing bases for $\mathbb{V}_1, \dots, \mathbb{V}_p, \mathbb{W}_1, \dots, \mathbb{W}_q$ and representing $X, Y$
as discrete functions by its representations in those basis, we get a way for computing $\mathcal{C}_k^l(X \otimes Y)$:

\begin{definition}[Base dependant tensor contraction]
    \normalfont{\cite{zhengFullyConnectedTensorNetwork2021}} Given $\mathcal{X}, \mathcal{Y}$ as the discrete function representations
    of $X, Y$ from earlier, we can write $\mathcal{C}_k^l(X \otimes Y)$ element-wise as:
    \begin{equation}
    \begin{align}\mathcal{C}_k^l & (i_1, \dots, i_{k-1}, i_{k+1}, \dots, i_p, j_1, \dots, j_{l-1}, j_{l+1}, \dots, j_q) \\  =& \sum_{s=1}^{N_k}
        \mathcal{X}(i_1, \cdots, i_{k-1}, s, i_{k+1}, \cdots, i_p) \mathcal{Y}(j_1, \cdots, j_{l-1}, s, j_{l+1}, \cdots, j_q)
\end{align}
\label{eq:contraction}
    \end{equation}
    \end{definition}

\begin{example}
    The tensor contraction for two tensors of order $2$ $X \in \mathbb{K}^{N_1 \times N_2}, Y \in \mathbb{K}^{N_2 \times N_3}$ over one edge of each tensor yields the matrix multiplication.
    Applying \refeq{eq:contraction} we get that
    $$(\mathcal{X} \times_2^2 \mathcal{Y}) (i_1, j_2) = \sum_{s=1}^{N_2} X(i_1, s) Y(s, j_2)$$
    And it is equivalent to the conventional matrix product. Visually, using the Penrsone notation we have:
    \begin{center}
    \hfill
\begin{minipage}{0.2\textwidth}
    \begin{tikzpicture}
    \node(X)[draw, shape=circle, minimum size=0.8cm] at (0,0) {$\mathcal{X}$};
    \node(Y)[draw, shape=circle, minimum size=0.8cm] at (2,0) {$\mathcal{Y}$};
    \node(i1) at (-1, 0) {$i_1$};
    \node(j2) at (3, 0) {$j_2$};
    \draw (X) -- (i1);
    \draw (Y) -- (j2);
    \draw (X) -- (Y) 
        node[above, pos=0.8] {$j_1$}
        node[above, pos=0.2] {$i_2$};
\end{tikzpicture}
\end{minipage}
\hfill
\begin{minipage}{0.225\textwidth}
\begin{tikzpicture}
    \node(1) at(-1.5, 1);
    \node(2) at(3, 1);
    \node(i) at(3, 0.5);
    \draw[-{Stealth[length=12pt, inset=2pt]}] (1) -- (2)
        node[above, pos=0.45] {$\mathcal{C}_{i_2}^{j_1}$};
\end{tikzpicture}
\end{minipage}
\hfill
\begin{minipage}{0.225\textwidth}
\begin{tikzpicture}
    \node(Z)[draw, shape=circle, minimum size=0.8cm] at (0,0) {$\mathcal{Z}$};
    \node(i1) at (-1, 0) {$i_1$};
    \node(j2) at (1, 0) {$j_2$};
    \draw (Z) -- (i1);
    \draw (Z) -- (j2);
\end{tikzpicture}
\end{minipage}

    \end{center}
\end{example}

\begin{example}
    Suppose that $X \in \mathbb{R}^{3 \times 3 \times 2}, Y \in \mathbb{R}^{2 \times 3 \times 3}$ with their representation
    in some given basis
    $$X = \left[ 
    \begin{pmatrix}
        1 & 0 & -1 \\
        0 & 2 & -2 \\
        -1 & 3 & 1
    \end{pmatrix} \quad
    \begin{pmatrix}
        1 & 1 & 0 \\
        2 & 0 & 2 \\
        0 & 3 & 0
    \end{pmatrix}
    \right] $$
  $$Y = \left[ 
    \begin{pmatrix}
        1 & 0 & -1 \\
        0 & 2 & -2 \\
    \end{pmatrix} \quad
    \begin{pmatrix}
        1 & 1 & 0 \\
        2 & 0 & 2 \\
    \end{pmatrix} \quad
    \begin{pmatrix}
        0 & 1 & 2 \\
        1 & 0 & 2
    \end{pmatrix}
    \right] $$

The contraction described in \figref{fig:tencon} would yield the tensor $\mathcal{Z} = \mathcal{X} \times_3^1 \mathcal{Y} \in \mathbb{R}^{3 \times 3 \times 3 \times 3}$ with
each element defined as:

$$\mathcal{Z}(i_1, i_2, j_2, j_3) = \sum_{s=1}^2 \mathcal{X}(i_1, i_2, s) \mathcal{Y}(s, j_2, j_3)$$

Therefore,

$$\mathcal{Z} = \begin{pmatrix}
\begin{pmatrix}
    1 & 2 & -3 \\
    3 & 1 & 2 \\
    1 & 1 & 4 \\
\end{pmatrix} &
\begin{pmatrix}
    0 & 2 & -2 \\
    2 & 0 & 2 \\
    1 & 0 & 2
\end{pmatrix} &
\begin{pmatrix}
    -1 & 0 & 1 \\
    -1 & -1 & 0 \\
    0 & -1 & -2
\end{pmatrix} \\
\begin{pmatrix}
    0 & 4 & -4 \\
    4 & 0 & 4 \\
    2 & 0 & 4
\end{pmatrix} &
\begin{pmatrix}
    2 & 0 & -2 \\
    2 & 2 & 0 \\
    0 & 2 & 4
\end{pmatrix} &
\begin{pmatrix}
    -2 & 4 & 6 \\
    2 & -2 & 4 \\
    2 & -2 & 0
\end{pmatrix} \\
\begin{pmatrix}
    -1 & 0 & 1 \\
    -1 & -1 & 0 \\
    0 & -1 & -2
\end{pmatrix} &
\begin{pmatrix}
    3 & 6 & -9 \\
    9 & 3 & 6 \\
    3 & 3 & 12
\end{pmatrix} &
\begin{pmatrix}
    1 & 0 & -1 \\
    1 & 1 & 0 \\
    0 & 1 & 2
\end{pmatrix} \\
\end{pmatrix}$$
\label{ex:bigmatrix}
\end{example}

\subsection *{Matrix tensor contraction}

Note that \refeq{eq:contraction} is almost in the form of matrix product. If we reshape accordingly
the tensors $\mathcal{X}$ and $\mathcal{Y}$, we can compute $\mathcal{C}_l^k(X \otimes Y)$ as a matrix product.
\begin{corollary}
    Let $X \in \mathbb{K}^{N_1 \times \cdots \times N_k \times \cdots \times N_n}, Y \in \mathbb{K}^{M_1 \times \cdots \times M_l \times \cdots \times M_m}$
    with $1 \leqslant k \leqslant n$, $1 \leqslant l \leqslant m$, $N_k = M_l$. The matrix product
$$ \mathcal{X}(\overline{i_1, \dots, i_{k-1}}, \overline{i_k, \dots, i_n}) \cdot \mathcal{Y}(\overline{j_1, \dots, j_{l-1}}, \overline{j_l, \dots, j_m}) $$
    results in a representation of $\mathcal{T} = \mathcal{X} \times_k^l \mathcal{Y}$
\end{corollary}

\begin{example}
    Following from \exref{ex:bigmatrix}, we would write $X \in \mathbb{R}^{9 \times 2}$ and $Y \in \mathbb{R}^{2 \times 9}$ as:

    $$X = \begin{pmatrix}
        1 & 1 \\ 0 & 1 \\ -1 & 0 \\ 0 & 2 \\ 2 & 0 \\ -2 & 2 \\ -1 & 0 \\ 3 & 3 \\ 1 & 0
        \end{pmatrix} \qquad
        Y = \begin{pmatrix}
            1 & 0 & -1 & 1 & 1 & 0 & 0 & 1 & 2 \\
            0 & 2 & -2 & 2 & 0 & 2 & 1 & 0 & 2
        \end{pmatrix}$$
        Therefore:
    $$Z = XY = \begin{pmatrix}
        1&	2&	-3&	 3&	 1&	2&	1&	 1&	 4\\
        0&	2&	-2&	 2&	 0&	2&	1&	 0&	 2\\
        -1&	0&	 1&	-1&	-1&	0&	0&	-1&	-2\\
        0&	4&	-4&	 4&	 0&	4&	2&	 0&	 4\\
        2&	0&	-2&	 2&	 2&	0&	0&	 2&	 4\\
        -2&	4&	-2&	 2&	-2&	4&	2&	-2&	 0\\
        -1&	0&	 1&	-1&	-1&	0&	0&	-1&	-2\\
        3&	6&	-9&	 9&	 3&	6&	3&	 3&	12\\
        1&	0&	-1&	 1&	 1&	0&	0&	 1&	 2\\
    \end{pmatrix}$$

    And we can now reshape $Z \in \mathbb{R}^{9 \times 9}$ as the tensor in $\mathbb{R}^{3\times 3\times 3\times 3}$
    as $$\mathcal{Z}(i_1, i_2, i_3, i_4) = Z(\overline{i_1, i_2}, \overline{i_3, i_4})$$
    which in fact, is identical to $\mathcal{Z}$ in \exref{ex:bigmatrix}
\end{example}

\begin{figure}
\centering
\tikz {
    \node(T)[draw, shape=circle, minimum size=0.8cm] at (0,0) {$\mathcal{T}$};
    \node(i) at (-1, 0) {$i_1$};
    \node(i2) at (-0.7, 0.8) {$i_2$};
    \node(ip) at (0, -1) {$\dots$};
    \node(iq) at (0, 1) {$\dots$};
    \node(in) at (-0.6, -0.8) {$i_n$};
    \node(ik) at (0.6, 0.5) {$i_k$};
    \node(ik1) at (1, -0.7) {$i_{k+1}$}
    \draw (T) -- (i);
    \draw (T) -- (i2);
    \draw (T) -- (in);
    \draw (T) -- (ip);
    \draw (T) -- (iq);
    \draw[--] (T) edge[in=15,out=-40,loop] ();
}
\caption{
    A tensor represented with the Penrose Notation with its indexes $i_k$ and $i_{k+1}$ connected as a loop
}
\label{fig:loop}
\end{figure}

There may be the case that when contracting a series of tensors, we might end up as what
we see as a loop in the Penrose Notation (see \figref{fig:loop}). Contracting over these two indexes we get the trace
of the tensor $\mathcal{T}$ respect the indices $i_k$ and $i_{k+1}$ and we denote it as $\Tr(T)$ (See \figref{fig:trace}):

\begin{definition}
    Given a tensor $\matchal{T} \in \mathbb{K}^{N_1 \times \cdots \times N_n}$ with $N_k = N_p$ with $1 \leqslant k, p \leqslant n$ and $k \neq p$,
    we define the trace of $\mathcal{T}$ respect the indices $k$ and $p$ as the tensor $\Tr_k^p(\mathcal{T})$ of order $n - 2$ with its entries defined as:
    $$\Tr_k^p(\mathcal{T}) (i_1, \dots, i_{k-1}, i_{k+1}, \dots, i_{p-1}, i_{p+1}, \dots, i_n) = \sum_{j = 1}^{N_k} \mathcal{T}(i_1, \dots, i_{k-1}, j, i_{k+1}, \dots, i_{p-1}, j, i_{p+1}, \dots, i_n)$$

\end{definition}

\begin{figure}[h]
\centering

    \begin{minipage}{0.2\textwidth}
        \begin{tikzpicture}
    \node(T)[draw, shape=circle, minimum size=0.8cm] at (0,0) {$\mathcal{T}$};
    \node(i) at (-1, 0) {$i_1$};
    \node(i2) at (-0.7, 0.8) {$i_2$};
    \node(i6) at (0, -1) {$i_6$};
    \node(i3) at (0, 1) {$i_3$};
    \node(i7) at (-0.6, -0.8) {$i_7$};
    \node(i4) at (0.6, 0.5) {$i_4$};
    \node(i5) at (1, -0.7) {$i_5$}
    \draw (T) -- (i);
    \draw (T) -- (i2);
    \draw (T) -- (i3);
    \draw (T) -- (i6);
    \draw (T) -- (i7);
    \draw[--] (T) edge[in=15,out=-40,loop] ();
\end{tikzpicture}
\end{minipage}
\begin{minipage}{0.225\textwidth}
    \begin{tikzpicture}
    \node(1) at(-1.5, -0.5);
    \node(2) at(1.5, -0.5);
    \draw[-{Stealth[length=12pt, inset=2pt]}] (1) -- (2)
        node[above, pos=0.45] {$\Tr(\mathcal{T})$};
\end{tikzpicture}
\end{minipage}
\begin{minipage}{0.2\textwidth}
    \begin{tikzpicture}
    \node(T)[draw, shape=circle, minimum size=0.8cm] at (0,0) {$\mathcal{T}$};
    \node(i) at (-1, 0) {$i_1$};
    \node(i2) at (-0.7, 0.8) {$i_2$};
    \node(i6) at (0, -1) {$i_6$};
    \node(i3) at (0, 1) {$i_3$};
    \node(i7) at (-0.6, -0.8) {$i_7$};
    \draw (T) -- (i);
    \draw (T) -- (i2);
    \draw (T) -- (i3);
    \draw (T) -- (i6);
    \draw (T) -- (i7);

    \end{tikzpicture}
\end{minipage}


\caption{
    Representation of the trace of a tensor using the Penrose notation
}
\label{fig:trace}
\end{figure}

Therefore, now one can ask that if we have an arbitrary graph $G$ represented with the Penrose Notation and it is connected,
we will be able to totally contract it to a single final tensor.


\section{Tensor Network States}

The concept of tensor networks originated from a physics background. Roger Penrose described how its
diagrammatic language could be used in various applications of physics \cite{rogerPenroseApplications}.

Later, in 1992, Steven R. White developed de Density Matrix Renormalization Group (DRMG) algorithm for
quantum lattice systems. It was considered the first successfull tensor network application \cite{whiteDensityMatrixFormulation1992}.

% Canviar aquest paràgraf
The commonly use terminology "tensor decomposition" is equivalent to "tensor network" to some extent. After
several years of progress accross different research fields, there is no significant distinction between tensor decomposition
and tensor networks. Tensor decomposition was employed primarly in signal processing fields \cite{wangTensorNetworksMeet2023}. Traditional
tensor decomposition models can be viewed
as basic kinds of tensor networks. In this thesis we will study some of the properties of different tensor decomposition methods,
and their effectivity.

We will define a mathematic definition of a tensor decomposition based on \cite{yeTensorNetworkRanks2019}, we will
see some common examples of tensor decompositions and we will define the 
tensor network structure space.

Informally, the way we construct a tensor network consists of picking a directed graph $G = (V, \bar{E})$, and for each vertex $i \in V$ we assign
a vector space $\mathbb{V}_i$ and for each edge $(i,j) \in \bar{E}$ we assign a vector space $\mathbb{E}_i$ to the tail of the edge and its 
dual covector space $\mathbb{E}_i^*$ to the head of the edge. We will also demand that the graph $G$ is connected.

More formally, let $\mathbb{V}_1, \dots, \mathbb{V}_d$ be vector spaces with $\dim{\mathbb{V}_i} = N_i, i = 1, \dots, d$. Let
$\mathbb{E}_1, \dots, \mathbb{E}_c$ be finite vector spaces with $\dim{\mathbb{E}_i} = R_i, i = 1, \dots, c$. For each $i \in V$ we associate
the tensor product space
$$\left( \bigotimes\nolimits_{j \in \IN (i)} \mathbb{E}_j \right) \otimes \mathbb{V}_i \otimes 
\left( \bigotimes\nolimits_{j \in \OUT (i)} \mathbb{E}_j^*  \right)$$

and a contraction map $\kappa_G$ defined by contracting factors in $\mathbb{E}_j$ with factors of $\mathbb{E}^*_j$

$$\kappa_G : \bigotimes\nolimits_{i=1}^d \left[ \left( \bigotimes\nolimits_{j \in \IN (i)} \mathbb{E}_j \right) \otimes \mathbb{V}_i \otimes 
\left( \bigotimes\nolimits_{j \in \OUT (i)} \mathbb{E}_j^*  \right)  \right] \rightarrow \bigotimes\nolimits_{i=1}^d \mathbb{V}_i$$

Note that we have given this shapes to the tensors that we fix onto each vertex
$i \in V$ because when we contract the whole graph, we will get a tensor of ${\mathbb{V}_1 \otimes \cdots \otimes \mathbb{V}_n}$.
Since every directed edge $(i,j)$ must point out of a vertex $i$ and point into a vetex $j$, each copy of $\mathbb{E}_j$ is paired with one
copy of $\mathbb{E}^*_$, so the contraction $\kappa_G$ is well defined and it results in a tensor of $\mathbb{V}_1 \otimes \cdots \otimes \mathbb{V}_n$
(See \figref{fig:tn_graph})

\begin{figure}[h]
    \centering

    \begin{minipage}{0.3\textwidth}
\begin{tikzpicture}[scale=0.85, arrows = {-Latex[length=8pt, inset=2pt]}]
    \node(1)[draw, shape=circle, minimum size=0.8cm] at (-3,1.5) {$1$};
    \node(2)[draw, shape=circle, minimum size=0.8cm] at (-3,-1.5) {$2$};
    \node(3)[draw, shape=circle, minimum size=0.8cm] at (0,0) {$3$};
    \node(4)[draw, shape=circle, minimum size=0.8cm] at (3,0) {$4$};
    \node(i1) at (-3, 2.5) {$\mathbb{V}_1$};
    \node(i2) at (-4, -1.5) {$\mathbb{V}_2$};
    \node(i3) at (0, 1) {$\mathbb{V}_3$};
    \node(i4) at (3, 1) {$\mathbb{V}_4$};


    \draw (1) -> (2)
    node[left, pos=0.8] {$\mathbb{E}_1$}
    node[left, pos=0.15] {$\mathbb{E}_1^*$};
    \draw (2) -> (3)
    node[above, pos=0.75] {$\mathbb{E}_2$}
    node[above, pos=0.1] {$\mathbb{E}_2^*$};
    \draw (3) -> (1)
    node[above, pos=0.75] {$\mathbb{E}_3$}
    node[above, pos=0.15] {$\mathbb{E}_3^*$};
    \draw (3) -> (4)
    node[above, pos=0.8] {$\mathbb{E}_4$}
    node[above, pos=0.15] {$\mathbb{E}_4^*$};
    \draw[--] (1) -- (i1);
    \draw[--] (2) -- (i2);
    \draw[--] (3) -- (i3);
    \draw[--] (4) -- (i4);
\end{tikzpicture}
\end{minipage}
\hfill
\begin{minipage}{0.28\textwidth}
    \begin{tikzpicture}
    \node(1) at(-1.5, -0.5);
    \node(2) at(2.5, -0.5);
    \draw[-{Stealth[length=12pt, inset=2pt]}] (1) -- (2)
        node[above, pos=0.45] {$\kappa_G(\mathbb{V}_1 \otimes \cdots \otimes \mathbb{V}_d)$};
\end{tikzpicture}
\end{minipage}
\begin{minipage}{0.25\textwidth}
    \begin{tikzpicture}
        \node(T)[draw, shape=circle, minimum size=0.8cm] at(0,0) {$T$};
   \node(i1) at (-1, 0) {$\mathbb{V}_1$};
    \node(i2) at (-0.4, 1) {$\mathbb{V}_2$};
    \node(i3) at (0.4, 1) {$\mathbb{V}_3$};
    \node(i4) at (1, 0) {$\mathbb{V}_4$};
    \draw (T) -- (i1);
    \draw (T) -- (i2);
    \draw (T) -- (i3);
    \draw (T) -- (i4);
    \end{tikzpicture}
\end{minipage}
\caption{
    Example of the correspondance between some directed graph $G$ and the vector spaces $\mathbb{V}_1, \dots, \mathbb{V}_d, \mathbb{E}_1, \dots, \mathbb{E}_c, \mathbb{E}_1^*,
    \dots, \mathbb{E}_c^*$ using the Penrose notation
}
\label{fig:tn_graph}
\end{figure}


\begin{definition}[Tensor network state]
    If a tensor $T \in \mathbb{V}_1 \otimes \cdots \otimes \mathbb{V}_n$ can be written as $T = \kappa_G(\mathcal{G}_1 \otimes \cdots \otimes \mathcal{G}_n)$ where
    $$\mathcal{G}_i \in \left( \bigotimes\nolimits_{j \in \IN (i)} \mathbb{E}_j \right) \otimes \mathbb{V}_i \otimes 
\left( \bigotimes\nolimits_{j \in \OUT (i)} \mathbb{E}_j^*  \right)$$
then we will say that $T$ is a \textbf{tensor network state} associated to the graph $G$ with cores $\mathcal{G}_i \in \mathbb{V}_i, i = 1, \dots, n$
\end{definition}

\begin{definition}[Tensor network]
    We will define all the resulting tensors that are possible by varying $\mathcal{G}_1, \dots, \mathcal{G}_n$ and then contracting through $\kappa_G$ as the tensor network
associated to $G$ and the vector spaces $\mathbb{V}_1, \dots, \mathbb{V}_n, \mathbb{E}_1, \dots, \mathbb{E}_c$ and we will write this set
as $\TNS(G; \mathbb{E}_1, \dots, \mathbb{E}_c, \mathbb{V}_1, \dots, \mathbb{V}_n)$, i.e
$$
\begin{align}
    \TNS(G; \mathbb{E}_1, \dots, \mathbb{E}_c, \mathbb{V}_1, \dots, \mathbb{V}_n) := & \Bigg\{ \kappa_G(\mathcal{G}_1 \otimes \cdots \otimes \mathcal{G}_n) \in \mathbb{V}_1 \otimes \cdots \otimes \mathbb{V}_n : \\
                                                                                     & \mathcal{G}_i \in \left( \bigotimes\nolimits_{j \in \IN (i)} \mathbb{E}_j \right) \otimes \mathbb{V}_i \otimes 
\left( \bigotimes\nolimits_{j \in \OUT (i)} \mathbb{E}_j^*  \right) \Bigg\}
\end{align}
$$
\end{definition}
Now, since all vector spaces are determined up to isomorphism by its dimension, when the vector spaces $\mathbb{E}_1, \dots, \mathbb{E}_c, \mathbb{V}_1, \dots, \mathbb{V}_n$ are unimportant,
we will write the tensor network as $\TNS(G; R_1, \dots, R_c, N_1, \dots, N_n)$. On the Penrose
Notation representation of the network, we will not write the vector spaces since they are
not important.

Also, since $n$ is equal to the number of vertices of $G$, we will sometimes write $\TNS(G; R)$ for a more compact notation.

\begin{definition}
    Given a tensor network state $\TNS(G, R)$, from the graph given by its Penrose Notation,
    we will call the edges with a dangling end \textit{free edges}, and the edges that connect two vertex \textit{contracted edges}
\end{definition}
For example, in \figref{fig:tn_graph}, the edges labeled as $\mathbb{V}_1, \mathbb{V}_2, \mathbb{V}_3, \mathbb{V}_4$ are 
free edges and the rest are contracted edges.

The following theorem gives us that if we make $R_1, \dots, R_c$ big enough, every tensor $T$ can be a state of $\TNS(G; R_1, \dots, R_c, N_1, \dots, N_n)$.
In fact, these values that guarantees that $T$ is an state are $R_1 = \dots R_c = \rank T$ where $\rank T$ is the \textit{traditional rank}
of a tensor. We will give now the definition of the traditional rank:

\begin{definition} 
    We say that a tensor has \textbf{traditional rank} $r$ and we write $\rank{T} = r$
    with $r \in \mathbb{N}$ being the minimum value such that we can write $T$ as
$$T= \sum_{p=1}^r v_p^{(1)} \otimes \cdots \otimes v_p^{(n)}$$
where $v_1^{(i)}, \dots, v_r^{(i)} \in \mathbb{V}_i, i = 1, \dots, n$
\end{definition}

We can immediatly see that $\rank{T} \leqslant \prod_{i=1}^n N_i$.
Unlike matrices, determining the rank of a tensor is an NP-hard problem. \cite{hillarMostTensorProblems2013}. Even
finding the maximum traditional rank of a tensor space (that means finding $\max_{T \in \mathbb{V}_1 \otimes \cdots \otimes \mathbb{V}_n} \rank{T}$) it still remains as an open problem.
We can find an slightly better upper bound for the tensor rank:
\begin{proposition}
    \begin{equation} \label{eq:rank-dimensional-bound}
        \rank{T} \leqslant \left\lfloor \frac{\prod_{i=1}^n N_i}{\sum_{i=1}^n N_i} \right\rfloor
    \end{equation}

\begin{proof}
    Let $r = \rank{T}$. We can write $T = \sum_{p=1}^r v_p^{(1)} \otimes \cdots \otimes v_p^{(n)}$. Now, each term of this
    sum has $\sum_{i=1}^n N_i$ adjustable parameters, since each $v_p^{(i)}$ is a vector of $\mathbb{V}_i$ with its dimension being $N_i$.
    So, in total we will have $r \sum_{i=1}^n N_i$ adjustable parameters in our decomposition. Since our tensor $T$ is completly
    determined by $\prod_{i=1}^n N_i$ parameters, we can impose that $r \sum_{i=1}^n N_i \leqslant \prod_{i=1}^n N_i$ 
\end{proof}
\end{proposition}

\begin{theorem} \normalfont{\cite{yeTensorNetworkRanks2019}} Let $T \in \mathbb{V}_1 \otimes \cdots \otimes \mathbb{V}_n$ and let $G$ be a connected graph with $n$ vertices and $c$ edges.
    There exists $R_1, \dots, R_c \in \mathbb{N}$ such that
    $$T \in \TNS(G; R_1, \dots, R_c, N_1, \dots, N_d)$$
    in fact, we can choose $R_1 = \dots = R_c = \rank{T}$
    \label{thm:finiterank}
\end{theorem}

\begin{proof} Let $r = \rank{T}$. Then there exist $v_1^{(i)}, \dots, v_r^{(i)} \in \mathbb{V}_i, i = 1, \dots, n$ such that
    $$T = \sum_{p=1}^r v_1^{(p)} \otimes \cdots \otimes v_n^{(p)}$$
    We take $R_1 = \dots = R_c = r$ we take for each $i = 1, \dots, n$
    $$\mathcal{G}_i = \sum_{p=1}^r \left( \bigotimes\nolimits_{j \in \IN(i)} e_p^{(j)} \right) \otimes v_p^{(i)} \otimes 
    \left( \bigotimes\nolimits_{j \in \OUT(i)} e_p^{(j)*} \right)$$

    Now observe that for each $i = 1, \dots, n$ there exists an unique $h$ such that whenever $j \in \IN(i)\cap \OUT(i)$,
    $e_p^{(j)}$ and $e_p^{(j)*}$ contract and give $\delta_{pq}$, therefore the summand vanishes except when $p = q$.
    This together with the assumption that $G$ is a connected graph implies that $\kappa_G(\mathcal{G}_1 \otimes \cdots \otimes \mathcal{G}_n)$ reduces
    to a sum of terms of the form $v_p^{(1)} \otimes \cdots \otimes v_p^{(d)}$ for $p = 1, \dots, r$, which is of course $T$
\end{proof}



\begin{definition} [Tensor G-rank]
    Given a graph $G$, we define the tensor rank respect to a $G$ as the mapping
    $$ \begin{align}
        \rank_G : \mathbb{V}_1 \otimes \cdots \otimes \mathbb{V}_n & \longrightarrow \mathcal{P}(\mathbb{N}^c)  \\ 
        T &\longmapsto \min {\{(R_1, \dots, R_c) \in \mathbb{N}^c : T \in \TNS(G; R_1, \dots, R_c, N_1, \dots, N_d)\}}
    \end{align}$$
 
    Where $\min(S)$ with $S \subset \mathbb{N}^c$ denotes the minimal elements of $S$. We treat $\mathbb{N}^c$ with its usual
    partial order:
    $$ (a_1, \dots, a_c) \leqslant (b_1, \dots, b_c) \Longleftrightarrow a_1 \leqslant b_1, a_2 \leqslant b_2, \dots, a_c \leqslant b_c$$
    So for example if $S = \{(3,4,5), (2,1,3), (1,3,2)\}$, then $\min(S) = \{(2,1,3), (1,3,2)\}$ Note that by \thmref{thm:finiterank} $\rank_G(T)$ is always a finite set.
\end{definition}

We will say that if $(r_1, \dots, r_c) \in \rank_G(T)$, we will say that $(r_1, \dots, r_c)$ is a $G$-rank of $T$. We will later see
that in some tensor network structures, for example in tree tensor networks, the $G$-rank of a tensor $T$ is unique



\section{Common Tensor network structures}

\begin{example}[Tensor Train decomposition]
    \normalfont{\cite{oseledetsTensorTrainDecomposition2011}} Let $\mathcal{T} \in \mathbb{R}^{N_1 \times \cdots \times N_n}$. A tensor
train decomposition or matrix product state of $\mathcal{T}$ are a set of $3$th-order tensors $\mathcal{G}_1,\mathcal{G}_2,\dots,\mathcal{G}_n$ with
    $\mathcal{G}_i \in \mathbb{R}^{R_{i-1} \times N_i \times R_i}$ and $R_0 = R_n = 1$ such that every element of $\mathcal{T}$ is written in the
    form
    \begin{equation} \label{eq:tt-contraction}
    \mathcal{T}(i_1,i_2,\dots,i_n) = \sum_{r_0, \dots, r_n}^{R_0, \dots, R_n} \mathcal{G}_1(r_0, i_1, r_1) \mathcal{G}_2 (r_1, i_2, r_2) \cdots \mathcal{G}_n(r_{n-1}, i_n, r_n)
\end{equation}
We denote $R_0, R_1, \dots, R_n$ as the ranks of the tensor train decomposition, or $TT$-ranks.
\end{example}

We can easily see that the tensor train decomposition (or TT) is obtained by our definition of a tensor network when $G$ is a path, also
the contraction of the whole network yields \refeq{eq:tt-contraction}


\begin{figure}[h]
    \centering
   \begin{minipage}{0.3\textwidth}
\begin{tikzpicture}[scale=0.85, arrows = {-Latex[length=8pt, inset=2pt]}]
    \node(1)[draw, shape=circle, minimum size=0.8cm] at (-4,0) {$\mathcal{G}_1$};
    \node(2)[draw, shape=circle, minimum size=0.8cm] at (-1,0) {$\mathcal{G}_2$};
    \node(3)[shape=circle, minimum size=0.8cm] at (2,0) {$\dots$};
    \node(4)[draw, shape=circle, minimum size=0.8cm] at (5,0) {$\mathcal{G}_n$};
    \node(i1) at (-4, 1) {$N_1$};
    \node(i2) at (-1, 1) {$N_2$};
    \node(i4) at (5, 1) {$N_n$};


    \draw (1) -> (2)
    node[above, pos=0.5] {$R_1$}
    \draw (2) -> (3)
    node[above, pos=0.5] {$R_2$}
    \draw (3) -> (4)
    node[above, pos=0.5] {$R_n$}
    \draw[--] (1) -- (i1);
    \draw[--] (2) -- (i2);
    \draw[--] (4) -- (i4);
\end{tikzpicture}
\end{minipage}
\hfill
\begin{minipage}{0.15\textwidth}
    \begin{tikzpicture}
    \node(1) at(-1, -0.5);
    \node(2) at(1, -0.5);
    \draw[-{Stealth[length=12pt, inset=2pt]}] (1) -- (2)
        node[above, pos=0.45] {$\kappa_G$};
\end{tikzpicture}
\end{minipage}
\begin{minipage}{0.25\textwidth}
    \begin{tikzpicture}
        \node(T)[draw, shape=circle, minimum size=0.8cm] at(0,0) {$T$};
   \node(i1) at (-1, 0) {$N_1$};
    \node(i2) at (-0.4, 1) {$N_2$};
    \node(i3) at (0.4, 1) {$\dots$};
    \node(i4) at (1, 0) {$N_n$};
    \draw (T) -- (i1);
    \draw (T) -- (i2);
    \draw (T) -- (i3);
    \draw (T) -- (i4);
    \end{tikzpicture}
\end{minipage}

    \caption{Tensor Train decomposition}
    \label{tt:scheme}
\end{figure}

\iffalse
\begin{theorem}
    \normalfont{\cite{oseledetsTensorTrainDecomposition2011}}
    Given the unfoldings $A_k = \mathcal{T}_{[1:k, k:n]}$, if we choose $\rank{A_k} = r_k$ then there exists a $TT$-decomposition
    with its ranks not higher than $r_k$
\end{theorem}
\fi


\begin{example}[Tensor Ring decomposition]
    \normalfont{\cite{zhaoTensorRingDecomposition2016}} Tensor ring decomposition (or TR) or also known a matrix product state with periodic boundary conditions, is obtained when $G$ is a cycle.
\end{example}
Tensor Ring decomposition is considered generalization of Tensor Train decomposition, it's contraction is the same as
\refeq{eq:tt-contraction} but removing the condition $R_0 = R_1 = 1$.

\begin{figure}[h]

    \centering
   \begin{minipage}{0\textwidth}
\begin{tikzpicture}[scale=0.85, arrows = {-Latex[length=8pt, inset=2pt]}]
    \node(1)[draw, shape=circle, minimum size=0.8cm] at (-1,2) {$\mathcal{G}_1$};
    \node(2)[draw, shape=circle, minimum size=0.8cm] at (1,2) {$\mathcal{G}_2$};
    \node(3)[draw, shape=circle, minimum size=0.8cm] at (2,0) {$\mathcal{G}_3$};
    \node(4)[draw, shape=circle, minimum size=0.8cm] at (1,-2) {$\mathcal{G}_4$};
    \node(5)[shape=circle, minimum size=0.8cm] at (-1,-2) {$\dots$};
    \node(n)[draw, shape=circle, minimum size=0.8cm] at (-2,0) {$\mathcal{G}_n$};

    \node(i1) at (-1.5, 3) {$N_1$};
    \node(i2) at (1.5, 3) {$N_2$};
    \node(i3) at (3.25, 0) {$N_3$};
    \node(i4) at (1.5, -3) {$N_4$};
    \node(in) at (-3.25, 0) {$N_n$};

    \draw (1) -> (2)
    node[above, pos=0.5] {$R_1$}
    \draw (2) -> (3)
    node[above right, pos=0.5] {$R_2$}
    \draw (3) -> (4)
    node[below right, pos=0.5] {$R_3$}
    \draw (4) -> (5)
    node[below, pos=0.5] {$R_4$}
    \draw (5) -> (n)
        node[right, pos=0.5] {$R_{n-1}$}
    \draw (n) -> (1)
    node[below right, pos=0.5] {$R_n$}

    \draw[--] (1) -- (i1);
    \draw[--] (2) -- (i2);
    \draw[--] (3) -- (i3);
    \draw[--] (4) -- (i4);
    \draw[--] (n) -- (in);
\end{tikzpicture}
\end{minipage}
\hfill
\begin{minipage}{0.235\textwidth}
    \begin{tikzpicture}
    \node(1) at(-1.5, -0.5);
    \node(2) at(1.5, -0.5);
    \draw[-{Stealth[length=12pt, inset=2pt]}] (1) -- (2)
        node[above, pos=0.45] {$\kappa_G$};
\end{tikzpicture}
\end{minipage}
\begin{minipage}{0.3\textwidth}
    \begin{tikzpicture}
        \node(T)[draw, shape=circle, minimum size=0.8cm] at(0,0) {$T$};
   \node(i1) at (-1, 0) {$N_1$};
    \node(i2) at (-0.4, 1) {$N_2$};
    \node(i3) at (0.4, 1) {$\dots$};
    \node(i4) at (1, 0) {$N_n$};
    \draw (T) -- (i1);
    \draw (T) -- (i2);
    \draw (T) -- (i3);
    \draw (T) -- (i4);
    \end{tikzpicture}
\end{minipage}

    \caption{Tensor Ring (TR) decomposition}
    \label{tr:schema}
\end{figure}

\begin{theorem} [Circular dimensional permutation invariance] Let $\mathcal{T} \in \mathbb{R}^{N_1 \times \cdots \times N_n}$ be
    a $n$th-order tensor with its corresponding tensor ring decomposition $\mathcal{T} = \mathcal{R}(\mathcal{U}^{(1)}, \mathcal{U}^{(2)}, \dots, \mathcal{U}^{(n)})$

\end{theorem}

\begin{example}[Fully connected tensor network]
The fully connected tensor network decomposition is obtenied when $G$ is a complete graph.
\end{example}

\begin{figure}[h]

    \centering
   \begin{minipage}{0\textwidth}
\begin{tikzpicture}[scale=0.75, arrows = {-Latex[length=8pt, inset=2pt]}]
    \node(1)[draw, shape=circle, minimum size=0.8cm] at (-1.5,3) {$\mathcal{G}_1$};
    \node(2)[draw, shape=circle, minimum size=0.8cm] at (1.5,3) {$\mathcal{G}_2$};
    \node(3)[draw, shape=circle, minimum size=0.8cm] at (3,0) {$\mathcal{G}_3$};
    \node(4)[draw, shape=circle, minimum size=0.8cm] at (1.5,-3) {$\mathcal{G}_4$};
    \node(5)[shape=circle, minimum size=0.8cm] at (-1.5,-3) {$\dots$};
    \node(n)[draw, shape=circle, minimum size=0.8cm] at (-3,0) {$\mathcal{G}_n$};

    \node(i1) at (-2, 4) {$N_1$};
    \node(i2) at (2, 4) {$N_2$};
    \node(i3) at (4.5, 0) {$N_3$};
    \node(i4) at (2, -4.5) {$N_4$};
    \node(in) at (-4.5, 0) {$N_n$};

    \draw (1) -> (2)
    \draw (2) -> (3)
    \draw (3) -> (4)
    \draw (4) -> (5)
    \draw (5) -> (n)
    \draw (n) -> (1)

    \draw (1) -> (3)
    \draw (1) -> (4)

    \draw (2) -> (n)
    \draw (2) -> (4)

    \draw (3) -> (n)

    \draw (1) -> (5)
    \draw (2) -> (5)
    \draw (3) -> (5)


    \draw[--] (1) -- (i1);
    \draw[--] (2) -- (i2);
    \draw[--] (3) -- (i3);
    \draw[--] (4) -- (i4);
    \draw[--] (n) -- (in);
\end{tikzpicture}
\end{minipage}
\hfill
\begin{minipage}{0.235\textwidth}
    \begin{tikzpicture}
    \node(1) at(-1.5, -0.5);
    \node(2) at(1.5, -0.5);
    \draw[-{Stealth[length=12pt, inset=2pt]}] (1) -- (2)
        node[above, pos=0.45] {$\kappa_G$};
\end{tikzpicture}
\end{minipage}
\begin{minipage}{0.3\textwidth}
    \begin{tikzpicture}
        \node(T)[draw, shape=circle, minimum size=0.8cm] at(0,0) {$T$};
   \node(i1) at (-1, 0) {$N_1$};
    \node(i2) at (-0.4, 1) {$N_2$};
    \node(i3) at (0.4, 1) {$\dots$};
    \node(i4) at (1, 0) {$N_n$};
    \draw (T) -- (i1);
    \draw (T) -- (i2);
    \draw (T) -- (i3);
    \draw (T) -- (i4);
    \end{tikzpicture}
\end{minipage}

    \caption{Fully connected tensor network decomposition (FCTN)}
    \label{tctn:schema}
\end{figure}


This theorem serves as an upper bound for our problem. But either way it is still very big. We can take a different approach:
we can fix $G, R_1, \dots, R_c$ and then find the cores of the tensor network that approximate better to our
objective tensor $T$. That means finding

\begin{equation}
\argmin_{\mathcal{G}_1, \dots, \mathcal{G}_n} \|T - \kappa_G (\mathcal{G}_1, \dots, \mathcal{G}_n) \|_F
\label{eq:minarg}
\end{equation}

Later, we will discuss some algorithms that gives us a solution to \refeq{eq:minarg}. 
\section{Tree Tensor Networks}

When $G$ is a tree, i.e when there are no cycles inside $G$, we say that $\TNS(G; R)$ is a tree tensor network.
This includes all tensor train networks. There are a lot of interesting results in tree tensor networks. One
very useful property is that given a tensor $T$, his $G$ rank when $G$ is a tree, is always unique.

\begin{figure}[h]

    \centering
\begin{tikzpicture}[scale=0.75, arrows = {-Latex[length=8pt, inset=2pt]}]
    \node(1)[draw, shape=circle, minimum size=0.8cm] at (0,4) {$\mathcal{G}_1$};
    \node(2)[draw, shape=circle, minimum size=0.8cm] at (-2,1) {$\mathcal{G}_2$};
    \node(3)[draw, shape=circle, minimum size=0.8cm] at (0,1) {$\mathcal{G}_3$};
    \node(4)[draw, shape=circle, minimum size=0.8cm] at (2,1) {$\mathcal{G}_4$};
    \node(5)[draw, shape=circle, minimum size=0.8cm] at (-5,-2) {$\mathcal{G}_5$};
    \node(6)[draw, shape=circle, minimum size=0.8cm] at (-3,-2) {$\mathcal{G}_6$};
    \node(7)[draw, shape=circle, minimum size=0.8cm] at (-1,-2) {$\mathcal{G}_7$};
    \node(8)[draw, shape=circle, minimum size=0.8cm] at (2,-2) {$\mathcal{G}_8$};

    \node(i1) at (-1.5, 4) {$N_1$};
    \node(i2) at (-3.5, 1) {$N_2$};
    \node(i3) at (0, -0.5) {$N_3$};
    \node(i4) at (3.5, 1) {$N_4$};
    \node(i5) at (-5, -3.5) {$N_5$};
    \node(i6) at (-3, -3.5) {$N_6$};
    \node(i7) at (-1, -3.5) {$N_7$};
    \node(i8) at (2, -3.5) {$N_8$};

    \draw (1) -> (2);
    \draw (1) -> (3);
    \draw (1) -> (4);
    
    \draw (2) -> (5);
    \draw (2) -> (6);
    \draw (2) -> (7);

    \draw (4) -> (8);

    \draw[--] (1) -- (i1);
    \draw[--] (2) -- (i2);
    \draw[--] (3) -- (i3);
    \draw[--] (4) -- (i4);
    \draw[--] (5) -- (i5);
    \draw[--] (6) -- (i6);
    \draw[--] (7) -- (i7);
    \draw[--] (8) -- (i8);
\end{tikzpicture}

    \caption{A tree tensor network}
    \label{tctn:schema}
\end{figure}









\section{The Alternating Least Squares algorithm}

\cite{malikSamplingBasedDecompositionAlgorithms2022a}
Suppose we want to only optimize the function \refeq{eq:minarg} for 
only one variable core $\mathcal{G}_m$ leaving the rest fixed. Then our problem would become
$$\argmin_{\mathcal{G}_m} \|T - \kappa_G(\mathcal{G}_1, \dots, \mathcal{G}_n)\|_F$$
Now, we could apply our contraction mapping $\kappa_G$ for all cores excluding $\mathcal{G}_m$ (\figref{fig:als_tr,fig:als_con}).
We will call this tensor $\mathcal{G}^{\neq m}$. Now, if we consider appropiate matricizations $T^{(m)}$, $G^{\neq m}$ and $G_m$
of $T$, $\mathcal{G}^{\neq m}$ and $\mathcal{G}_i$ respectively, evaluating the whole tensor network
is equal to computing the product $G^{\neq m} G_m$, so our problem is equivalent to solve the 
following linear least squares problem:
\begin{equation}
    \argmin_{G_m} \| G^{\neq m} G_m - T^{(m)}\|_2
\label{eq:als}
\end{equation}
Let $x^{(i)}$ be the $i$-th column of $G_m$ and $y^{(i)}$ the $i$-th
column of $T^{(m)}$. Solving \ref{eq:als} means solving for each $i$ 
\begin{equation}
    \argmin_{x^{(i)}} \| G^{\neq m} x^{(i)} - y^{(i)} \|_2
\end{equation}

Since we can't assure that there exists an exact solution to $G^{\neq m} x^{(i)} = y^{(i)}$,
we can use the solution to the normal equation $(G^{\neq m})^T G^{\neq m} x^{(i)} = (G^{\neq m})^T y^{(i)}$.

Now we can iteratively change the varying core tensor $\mathcal{G}_i$ until the contraction of the whole tensor network
is $T$ with some fixed error $\epsilon$:

\begin{algorithm}
    \caption{Tensor Network ALS}

    \hspace*{\algorithmicindent} \textbf{Input}: A tensor $T \in \mathbb{K}^{N_1 \times \cdots \times N_n}$ and some fixed error $\epsilon$ \\
    \hspace*{\algorithmicindent} \textbf{Output}: Core tensors $\mathcal{G}_1, \dots, \mathcal{G}_n$ 

    \begin{algorithmic}[1]
        \State Initialize tensors $\mathcal{G}_1, \dots, \mathcal{G}_n$
        \While{$\|T - \kappa_G (\mathcal{G}_1, \dots, \mathcal{G}_n)\|_F > \epsilon$ }
            \For {$k = 1, \dots, n$}
            \State $\mathcal{G}_m \leftarrow \argmin_{G_m} \|G^{\neq m} G_m - T^{(m)}\|_2$
            \EndFor
        \EndWhile
        \State \Return $\mathcal{G}_1, \dots, \mathcal{G}_n$

    \end{algorithmic}

\end{algorithm}


\begin{figure}[h]

    \centering
   \begin{minipage}{0\textwidth}
       \begin{tikzpicture}[scale=0.85]
    \node(1)[draw, shape=circle, minimum size=0.8cm] at (-1,2) {$\mathcal{G}_2$};
    \node(2)[draw, shape=circle, fill=lightgray, minimum size=0.8cm] at (1,2) {$\mathcal{G}_3$};
    \node(3)[draw, shape=circle, minimum size=0.8cm] at (2,0) {$\mathcal{G}_4$};
    \node(4)[draw, shape=circle, minimum size=0.8cm] at (1,-2) {$\mathcal{G}_5$};
    \node(5)[draw, shape=circle, minimum size=0.8cm] at (-1,-2) {$\mathcal{G}_6$};
    \node(n)[draw, shape=circle, minimum size=0.8cm] at (-2,0) {$\mathcal{G}_1$};

\node(I) at(2.5, 4) {$\kappa_G(\mathcal{G}_1, \dots, \mathcal{G}_n)$}

    \node(i1) at (-1.5, 3) {$N_2$};
    \node(i2) at (1.5, 3) {$N_3$};
    \node(i3) at (3.25, 0) {$N_4$};
    \node(i4) at (1.5, -3) {$N_5$};
    \node(i5) at (-1.5, -3) {$N_6$};
    \node(in) at (-3.25, 0) {$N_1$};

    \draw (1) -> (2)
    node[above, pos=0.5] {$R_2$}
    \draw (2) -> (3)
    node[above right, pos=0.5] {$R_3$}
    \draw (3) -> (4)
    node[below right, pos=0.5] {$R_4$}
    \draw (4) -> (5)
    node[below, pos=0.5] {$R_5$}
    \draw (5) -> (n)
        node[right, pos=0.5] {$R_6$}
    \draw (n) -> (1)
    node[below right, pos=0.5] {$R_1$}

    \draw[--] (1) -- (i1);
    \draw[--] (2) -- (i2);
    \draw[--] (3) -- (i3);
    \draw[--] (4) -- (i4);
    \draw[--] (5) -- (i5);
    \draw[--] (n) -- (in);
\end{tikzpicture}
\end{minipage}
\hfill
\begin{minipage}{0.235\textwidth}
    \begin{tikzpicture}
        \node(a) at (-2, 0);
    \node(1) at(-1.5, 0);
    \node(2) at(1.25, 0);
    \draw[-{Stealth[length=12pt, inset=2pt]}] (1) -- (2);
\end{tikzpicture}
\end{minipage}
\begin{minipage}{0.35\textwidth}
    \begin{tikzpicture}[scale=0.85]
    \node(3)[draw, shape=circle, minimum size=0.8cm] at (0,4) {$\mathcal{G}_4$};
    \node(4)[draw, shape=circle, minimum size=0.8cm] at (0,2) {$\mathcal{G}_5$};
    \node(5)[draw, shape=circle, minimum size=0.8cm] at (0, 0) {$\mathcal{G}_6$};
    \node(0)[draw, shape=circle, minimum size=0.8cm] at (0, -2) {$\mathcal{G}_1$};
    \node(1)[draw, shape=circle, minimum size=0.8cm] at (0, -4) {$\mathcal{G}_2$};

    \node(ii1) at (0, 5) {};
    \node(i3) at (-1.5, 4) {$N_4$};
    \node(i4) at (-1.5, 2) {$N_5$};
    \node(i5) at (-1.5, 0) {$N_6$};
    \node(i0) at (-1.5, -2) {$N_1$};
    \node(i1) at (-1.5, -4) {$N_2$};
    \node(ii2) at (0, -5) {};

    \node(I) at (2, 4) {$\mathcal{G}^{\neq 2}$};

    \draw (3) -> (4)
    node[right, pos=0.5] {$R_4$}
    \draw (4) -> (5)
    node[right, pos=0.5] {$R_5$}
    \draw (5) -> (0)
    node[right, pos=0.5] {$R_6$}
    \draw (0) -> (1)
    node[right, pos=0.5] {$R_1$}
    \draw[--] (1) -- (ii2)
    node[right] {$R_3$};
    \draw[--] (ii1) -- (3)
    node[above right, pos=1] {$R_2$};



    \draw[--] (1) -- (i1);
    \draw[--] (3) -- (i3);
    \draw[--] (4) -- (i4);
    \draw[--] (5) -- (i5);
    \draw[--] (0) -- (i0);
    \end{tikzpicture}
\end{minipage}

\caption{The representation of $\mathcal{G}^{\neq m}$ on the TR decomposition, with $m = 2$ }
    \label{fig:als_tr}
\end{figure}
%
%
%
%
%
%
%
%
\begin{figure}[b]
    \centering
   \begin{minipage}{0\textwidth}
       \begin{tikzpicture}[scale=0.8]
    \node(1)[draw, shape=circle, minimum size=0.8cm] at (0,4) {$\mathcal{G}_1$};
    \node(2)[draw, shape=circle, minimum size=0.8cm] at (3.5,4) {$\mathcal{G}_2$};
    \node(3)[draw, shape=circle, minimum size=0.8cm] at (-2,2) {$\mathcal{G}_3$};
    \node(4)[draw, fill=lightgray, shape=circle, minimum size=0.8cm] at (1.75,2) {$\mathcal{G}_4$};
    \node(5)[draw, shape=circle, minimum size=0.8cm] at (5,2) {$\mathcal{G}_5$};
    \node(6)[draw, shape=circle, minimum size=0.8cm] at (0,0) {$\mathcal{G}_6$};
    \node(7)[draw, shape=circle, minimum size=0.8cm] at (3.5,0) {$\mathcal{G}_7$};


    \node(I) at(1.75, 6.5) {$\kappa_G(\mathcal{G}_1, \dots, \mathcal{G}_n)$}

    \node(i1) at (0, 5.25) {$N_1$};
    \node(i2) at (3.5, 5.25) {$N_2$};
    \node(i3) at (-2, 3.25) {$N_3$};
    \node(i4) at (0.75, 2.75) {$N_4$};
    \node(i5) at (5, 3.25) {$N_5$};
    \node(i6) at (0, -1.25) {$N_6$};
    \node(i7) at (3.5, -1.25) {$N_7$};

    \draw[--] (1) -- (i1);
    \draw[--] (2) -- (i2);
    \draw[--] (3) -- (i3);
    \draw[--] (4) -- (i4);
    \draw[--] (5) -- (i5);
    \draw[--] (6) -- (i6);
    \draw[--] (7) -- (i7);

    \draw (1) -> (2)
    node[above, pos=0.5] {$R_1$};
    \draw (2) -> (5)
    node[below left, pos=0.5] {$R_2$};
    \draw (4) -> (2)
    node[above left, pos=0.5] {$R_3$};
    \draw (3) -> (6)
    node[below left, pos=0.5] {$R_4$};
    \draw (6) -> (4)
    node[above left, pos=0.5] {$R_5$};
    \draw (4) -> (7)
    node[left, pos=0.5] {$R_6$};
    \draw (7) -> (5)
    node[above left, pos=0.5] {$R_7$};
\end{tikzpicture}
\end{minipage}
\hfill
\begin{minipage}{0.235\textwidth}
    \begin{tikzpicture}
        \node(a) at (-2, 1);
    \node(1) at(-1.5, 0);
    \node(2) at(1.25, 0);
    \draw[-{Stealth[length=12pt, inset=2pt]}] (1) -- (2);
\end{tikzpicture}
\end{minipage}
\begin{minipage}{0.35\textwidth}
    \begin{tikzpicture}[scale=0.8]
    
        \node(I) at(1.75, 6.5) {$\mathcal{G}^{\neq 4}$}

    \node(1)[draw, shape=circle, minimum size=0.8cm] at (0,4) {$\mathcal{G}_1$};
    \node(2)[draw, shape=circle, minimum size=0.8cm] at (3.5,4) {$\mathcal{G}_2$};
    \node(3)[draw, shape=circle, minimum size=0.8cm] at (-2,2) {$\mathcal{G}_3$};
    \node(4)[shape=circle, minimum size=0.8cm] at (1.75,2);
    \node(5)[draw, shape=circle, minimum size=0.8cm] at (5,2) {$\mathcal{G}_5$};
    \node(6)[draw, shape=circle, minimum size=0.8cm] at (0,0) {$\mathcal{G}_6$};
    \node(7)[draw, shape=circle, minimum size=0.8cm] at (3.5,0) {$\mathcal{G}_7$};

    \node(i1) at (0, 5.25) {$N_1$};
    \node(i2) at (3.5, 5.25) {$N_2$};
    \node(i3) at (-2, 3.25) {$N_3$};
    \node(i5) at (5, 3.25) {$N_5$};
    \node(i6) at (0, -1.25) {$N_6$};
    \node(i7) at (3.5, -1.25) {$N_7$};

    \draw[--] (1) -- (i1);
    \draw[--] (2) -- (i2);
    \draw[--] (3) -- (i3);
    \draw[--] (5) -- (i5);
    \draw[--] (6) -- (i6);
    \draw[--] (7) -- (i7);

    \draw (1) -> (2)
    node[above, pos=0.5] {$R_1$};
    \draw (2) -> (5)
    node[below left, pos=0.5] {$R_2$};
    \draw (4) -> (2)
    node[above left, pos=0.5] {$R_3$};
    \draw (3) -> (6)
    node[below left, pos=0.5] {$R_4$};
    \draw (6) -> (4)
    node[above left, pos=0.5] {$R_5$};
    \draw (4) -> (7)
    node[left, pos=0.5] {$R_6$};
    \draw (7) -> (5)
    node[above left, pos=0.5] {$R_7$};

    \end{tikzpicture}
\end{minipage}

\caption{The representation of $\mathcal{G}^{\neq m}$ on the TR decomposition, with $m = 2$ }
    \label{fig:als_con}
\end{figure}

\section{Contracting the tensor network}

The ALS algorithm requires us to contract the network when we update the cores and
also when we do a contraction for checking for the stopping condition. The order in which we 
select to contract the cores affects directly the computational cost of the contraction. 

For example, consider the following tensor train network:

\begin{center}
    \begin{tikzpicture}

    \node(1)[draw, shape=circle, minimum size=0.8cm] at (0,0) {$X$};
    \node(2)[draw, shape=circle, minimum size=0.8cm] at (3,0) {$Y$};
    \node(3)[draw, shape=circle, minimum size=0.8cm] at (6,0) {$Z$};
    \node(4)[draw, shape=circle, minimum size=0.8cm] at (9,0) {$T$};

    \node(i1) at (0, 1) {$N_1$};
    \node(i2) at (3, 1) {$N_2$};
    \node(i3) at (6, 1) {$N_3$};
    \node(i4) at (9, 1) {$N_4$};

    \draw[--] (1) -- (i1);
    \draw[--] (2) -- (i2);
    \draw[--] (3) -- (i3);
    \draw[--] (4) -- (i4);

    \draw (1) -> (2)
    node[above, pos=0.5] {$R_1$};
    \draw (2) -> (3)
    node[above, pos=0.5] {$R_2$};
    \draw (3) -> (4)
    node[above, pos=0.5] {$R_3$};
    \end{tikzpicture}
\end{center}

The contraction of the whole network is given by

$$T_{ijkl} = \sum_{p,q,r}^{R_1, R_2, R_3} X_{ip} Y_{pjq} Z_{qkr} T_{rl}$$

We could contract first $R_1$, then $R_2$ and finally $R_3$. That means computing:
$$A_{ijq} = \sum_{p}^{R_1} X_{ip} Y_{pjq} \qquad B_{ijkr} = \sum_{q}^{R_2} A_{ijq} Z_{qkr} \qquad T_{ijkl} = \sum_{r}^{R_3} B_{ijkr} T_{rl}$$

For this computation, we would need $R_1 R_2 N_1 N_2 + R_2 R_3 N_1 N_2 N_3 + R_3 N_1 N_2 N_3 N_4$ products, and we would need to store
first $A$, then $B$ and then $T$

If we consider a different ordering for the contraction, for example first $R_1$, then $R_3$ and finally $R_2$ we get a different computation cost.
We would need to compute the matrices

$$A'_{ijq} = \sum_{p}^{R_1} X_{ip} Y_{pjq} \qquad B'_{qkl} = \sum_{r}^{R_3} Z_{qkr} T_{rl} \qquad T_{ijkl} = \sum_{q}^{R_2} A'_{ijq} B'_{qkl}$$

The result of the contraction is the same, but now instead the number of computations is
changed to $R_1 R_2 N_1 N_2 + R_3 N_1 N_3 N_4 + R_2 N_1 N_2 N_3 N_4$. 


\chapter{Topological search on the tensor network structure space}

By now, we know that for any tensor $T$, if we choose a graph $G = (V, E)$ we can choose 
$R_1, \dots, R_c \leqslant \rank{T} \leqslant \left\lfloor \frac{\prod_{i=1}^n N_i}{\sum_{i=1}^n N_i} \right\rfloor$
such that $T \in \TNS(G, R)$. We can also find thanks to the ALS algorithm the cores $\mathcal{G}_1, \dots, \mathcal{G}_n$ of 
$\TNS(G, R)$. The only thing that remains is how we pick an optimal $G$ for compressing $T$ with a fixed error $\epsilon$, since we know that
the optimal graph for a tensor network state of a given tensor $T$ depends on the underlying data of $T$ itself.

We will aim to search the tensor network that gives the best compression ratio for $T$, for that, we will define the size 
of a tensor network:

\begin{definition}
    \cite{guoTensorNetworkStructure2025}
    Given a tensor network state $\TNS(G, R)$, we define $$\size(G, R) = \sum_{i=1}^c \size\mathcal{G}_i$$
\end{definition}

We would want to solve

$$\arg\min_{G, R} \;\size(G, R) \quad s.t \quad \|\kappa_G(\mathcal{G}_1, \dots, \mathcal{G}_c) - T\|_F \leqslant \epsilon \|T\|_F$$




Let $\mathcal{X} \in \mathbb{R}^{N_1 \times \cdots \times N_n}$ an $n$-order tensor. Let the following problem:

$$ \min_{r \in \mathbb{K}_{K_N} } \phi(K_N, r) \quad s.t\; \mathcal{X} \in \TNS{(K_N, r)} $$ 

\section{Counting tensor network structures}

- Fer alternating local enumeration


% \cite{liPermutationSearchTensor2022}



% TODO:
% Suposo que començar per dir quines parts del graf caldria tallar maybe???
% Fer més representacions gràfiques de segons quina demostració com més clar quedi tot millor
% \begin{itemize}
% \item Descriure $G$-ranks
% \item Algorismes per aproximar TNS per $G$-ranks propers i mínims si es pot fer
% \item Algun algorisme per trobar heuristicament els $G$-ranks adequats? (suposo q depen de compressió ratio i l'error relatiu)
% \item Com podem trobar un $G$ adequat?
% \item Estratègies per contraure tensors més ràpidament? (DRMG?)
% \item Algorismes, part pràctica en C/C++
% \item Fer moltes gràfiques
% \item Fer aplicacions per machine learning, etc.
% \item Fixar la mathematical subject classification
% \end{itemize}












\chapter{Conclusions}

TODO

Fent servir un s\'{\i}mil geom\`etrico-cartogr\`afic, aquesta mem\`oria constitueix un mapa a escala planet\`aria de la demostraci\'o de la conjectura feble de Goldbach presentada per Helfgott i un mapa a escala continental de la verificaci\'o num\`erica d'aquesta. Estudis posteriors i m\'es profunds haurien de permetre elaborar mapes de menor escala.

La naturalesa dels nombres primers ens ha portat per molts racons diferents de les Matem\`atiques; en no imposar-nos restriccions en la forma de pensar, hem pogut gaudir del viatge i assolir els objectius que ens vam plantejar a l'inici del projecte i anar m\'es enll\`a, sobretot en el camp de la computaci\'o i la manipulaci\'o de grans volums de dades num\`eriques.

Una gran part dels coneixements b\`asics que hem hagut de fer servir han estat treballats en les assignatures de M\`etodes anal\'{\i}tics en teoria de nombres i d'An\`alisi harm\`onica i teoria del senyal, que s\'on optatives de quart curs del Grau de Ma\-te\-m\`a\-ti\-ques. Altres els hem hagut d'aprendre durant el desenvolupant del projecte. S'ha realitzat una tasca de recerca bibliogr\`afica important, consultant recursos antics i moderns, tant en format digital com en format paper.

\normalfont


\newpage

\addcontentsline{toc}{chapter}{Bibliography}
\printbibliography

\appendix
\chapter{Chapter 1}

%\section{Tensor contractions}

%\begin{definition}[Tensor contraction]
%    (https://math.stackexchange.com/questions/1792230/coordinate-free-notation-for-tensor-contraction)
%\end{definition}

%Let $T \in \mathbb{V}_1 \otimes \cdots \otimes \mathbb{V}_k \otimes \cdots \otimes \mathbb{V}_p \otimes \mathbb{W}_1^* \otimes \cdots
%\otimes \mathbb{W}_l \otimes \cdots \otimes \mathbb{W}_q$

 
%\chapter{Experimental results}
\end{document} 

